% SHARD-ID: AISM-00A-DEFINITIONS
% SHARD-TITLE: The vocabulary
% SHARD-SUMMARY: Frames the generated definitions layer: every term the current proof strategy needs is stated once, typeset from its canonical definitions/ shard by scripts/gen-report-defs.py.
% SHARD-SUMMARY: Explains the two rules that govern the layer -- typeset statement first with the byte-verbatim source demoted to a second check, and scope restricted to the Route-F dependency closure -- plus the macro-translation table that makes a quoted source compile.
% SHARD-SUMMARY: Gives the rigour-ladder legend for definition kinds and statuses, the signed/stochastic reading guidance and the def:<slug> crosslink scheme, and records the honesty boundary: a locked cited definition certifies provenance, never truth, and a definition is never proved.
% SHARD-KEYWORDS: definitions, vocabulary, generated projection, macro-translation table, source check, proof-strategy scope, rigour ladder, cited, consensus, original, crosslinks, signed picture, stochastic picture
\section{The vocabulary}\label{sec:definitions}

\subsection*{What this section is, and why it is generated}

Every later section names its terms by \texttt{def-} identifier and never restates them.
This section is where those identifiers resolve. It reproduces the project's Layer-0
definitions database --- one shard per term under \texttt{definitions/} --- as typeset
definitions, in a deliberate reading order, each carrying its own provenance.

The material below is \emph{generated}, and that is a load-bearing design decision rather
than a convenience. The governing law of this repository is that each term is defined
exactly once (\texttt{CLAUDE.md}~L2): a second, hand-maintained statement of a definition
inside the paper is drift, and drift here is fatal, because a lemma proved against one
reading of $\negmass$ or of exposedness and quoted against another is a wrong result that
no gate would catch. So the canonical shard stays the single source of truth and the pages
below are a deterministic projection of it, produced by
\texttt{scripts/gen-report-defs.py} into \texttt{report/generated/defs/} in the same way
that \texttt{argument/INDEX.md} and \texttt{argument/DAG.md} are produced from the registry
shards. The generator's \texttt{--check} mode fails the local gate whenever the committed
projection differs from a fresh render, so a definition edited in \texttt{definitions/} and
not re-rendered here cannot be committed. Nothing in this section may be edited in place:
corrections belong in the named shard, after which the projection is regenerated.

\subsection*{The statement is the definition; the source text is the second check}

Each entry is a typeset \texttt{definition} environment, and that environment is the
definition. Underneath it sits a small provenance block; underneath that, where the shard
carries one, a \textsc{source check (byte-verbatim)} block. The order is deliberate. A
definition transcribed from the literature is only usable if it can be \emph{read}, and a
block of detokenized monospace is not a readable definition; but the byte-verbatim text is
what the repository's \texttt{cited} criterion actually locks (\texttt{L1}), so it cannot
simply be dropped. It is therefore kept and demoted: the small print at the end, against
which the typeset statement above may be checked, rather than the statement itself.

Where a shard carries a harmonised \textbf{Statement} section, that section is what is
typeset --- it exists in the shard precisely to be the readable form, and it is where a
transcription's known corrections are recorded. Where the shard's only content is the
pinned source's own \TeX{}, the generator typesets \emph{that}, through an explicit
macro-translation table. The table exists because a quoted source is written in that
source's private macro dialect ($\backslash$\texttt{eps}, $\backslash$\texttt{calA},
$\backslash$\texttt{Co}, \dots), which this document's preamble does not define. Every row
of the table maps a source macro to the expansion \emph{the source itself gives it}, taken
from that source's own preamble and recorded with the line number that defines it; the
table is reproduced in full in \texttt{report/generated/defs/MANIFEST.md} for audit, together
with its two declared glyph-only deviations. Translation changes spelling, never meaning.
Each definition typeset this way says so in its provenance block, names the table version,
and lists the macros that fired.

The honesty boundary is mechanical, not editorial. A translated fragment is validated before
it is emitted --- every surviving control sequence must be one this document defines, braces
and math delimiters must balance, no alignment or script character may escape math mode ---
and a fragment that fails is \emph{not} typeset. It is replaced on the page by a loud
\textbf{NOT TYPESET} block naming the reason, its bytes are left to the source check below,
and it is listed in \texttt{report/generated/defs/MANIFEST.md}. Partial success is allowed:
what translates is typeset, what does not is flagged, and nothing in between is guessed at.
As of this render there are no such flags. The byte-verbatim blocks themselves are set with
\TeX's detokenizer, which prints a control word followed by one space; apart from that
spacing artefact, and from the line breaks the quote block inserts to keep long source lines
inside the text block, what is printed is the source's own bytes.

\subsection*{Reading the kind and status of a definition}

Each definition is followed by a small provenance block naming its \emph{kind}, its
\emph{status}, its source, its locus, and --- for literature transcriptions --- the
\textsc{sha256} prefix of the pinned source file. The vocabulary of that block is the
repository's rigour ladder applied to definitions, and it is not the same ladder that
grades results.

\begin{itemize}
\item \texttt{cited} --- transcribed from the literature and byte-matched to a local source
  under \texttt{refs/} at the recorded locus and hash. The gate
  \texttt{scripts/check-defs.py} recomputes that hash. What this certifies is
  \emph{provenance}: that the text below is what the source says. It certifies nothing
  about whether the source's claim is true, and several of the cited shards here carry
  explicit scope notes recording exactly which neighbouring assertions of the source are
  \emph{not} part of the definition.
\item \texttt{consensus} --- a formulation this project adopted and agreed on, with no
  single external source to byte-match against.
\item \texttt{original} --- a notion introduced here, typically to name a configuration or
  a hypothesis package that the exploration kept re-inlining.
\end{itemize}

\noindent The status is orthogonal. \texttt{locked} means usable downstream: byte-verified
for a \texttt{cited} shard, signed off for a \texttt{consensus} or \texttt{original} one.
\texttt{draft} means the shard is pinned but not yet gated --- fine as working vocabulary,
and flagged whenever a validated result leans on it. A definition is never
\texttt{proved}: it carries no truth claim, so there is nothing in it to prove. When a
statement about a defined object needs to be true, it is a registry result with its own
status, not a clause smuggled into the vocabulary.

\subsection*{The two pictures, and which one a term lives in}

The single most reliable way to misread this project is to forget which picture a symbol
belongs to. The stochastic picture holds genuinely positive row-stochastic maps that are
only approximately idempotent (Definition~\ref{def:stochastic},
Definition~\ref{def:almost-idempotent}). The signed picture holds exactly idempotent
signed affine retractions whose failure of positivity is quarantined in a single scalar
(Definition~\ref{def:signed-idempotent}, Definition~\ref{def:negative-mass}); everything
built on row geometry --- exposedness, the visible set, height, invisible mass --- is
stated there. The licence to cross between the pictures, with explicit constants, is
\texttt{lem-classical-equiv} (\S\ref{sec:classical-equiv}), and it is the only such
licence.

Two symbol collisions are deliberate in the sources and are preserved rather than silently
renamed. The letter $\negmass$ denotes the negative mass of a signed idempotent in the
classical layer and the defect of a $\delta$-projection in the approximate-$C^{*}$ layer
(Definition~\ref{def:delta-projection}); these are different quantities, and
\texttt{CONVENTIONS.md} is the registry that keeps them apart. Likewise $\ecs$ is the
ambient approximate-$C^{*}$ parameter, never the retract defect of
Definition~\ref{def:positive-approximate-retract}. When a bound is quoted, the picture and
the parameter must be named with it.

\subsection*{Which definitions appear here}

This section renders the vocabulary of the \emph{current proof strategy}, not the whole
Layer-0 database. The strategy is the Route-F landing chain, and it is not re-described
here: the generator imports \texttt{scripts/gen-report-dag.py} and reuses the very same
subgraph computation that draws the atlas of \S\ref{sec:argument-dag}, so the map and the
vocabulary cannot drift apart. A definition is rendered when it is imported by a registry
result in that subgraph, or by a registry result reproduced in this document, or when it is
reached from one of those by following the \texttt{[[def-\dots]]} crosslinks that the
rendered statements themselves use. Anything outside that closure is not printed, and a
generated paragraph at the end of the section says how many such definitions there are and
names them; they remain canonical in \texttt{definitions/} and in
\texttt{definitions/INDEX.md}, they are simply not part of what this document needs. A
crosslink pointing at one of them degrades to a bare identifier rather than to a link this
document could not honour.

The three groups below are computed from that set, not curated. The first is the connected
component of the crosslink graph that the classical signed and stochastic vocabulary forms:
those definitions genuinely define each other, and they are ordered so that a term appears
after everything its statement uses, with mutually recursive clusters
(negative mass, signed idempotent, stochastic idempotent; exposedness and the visible set)
kept together. The second group is the approximate-$C^{*}$ vocabulary transcribed from the
pinned source, which is the language of the tiers reproduced in
\S\S\ref{sec:compcb-amplification-naturality}--\ref{sec:extcb}. The third is the
project-internal packaging built over that language: the datum bundles that fix the
hypotheses of a tier, and the notation the proofs of that tier share. A group with no
member in scope would not be printed at all.

One honest asymmetry survives the scoping. Being \emph{needed by} the strategy is a much
weaker condition than being \emph{proved on} it: much of the classical group is the
vocabulary of the signed-geometry trunk, whose results are ancestors of the target but are
recorded in \texttt{argument/} and whitelisted in \texttt{report/UNWIRED.md} rather than
reproduced here. Those definitions belong in this section --- the target
\texttt{op-classical} of \S\ref{sec:overview} is stated in their terms --- but the reader
should not read their presence as a claim that the surrounding argument is in this
document. Each definition's pointer line says how many registry results import it and which
of those, if any, appear here; a term with no reproduced consumer says so instead of
leaving the reader to guess.

\subsection*{Citing a definition}

Every definition carries \texttt{\textbackslash label\{def:}\textit{slug}\texttt{\}} and a matching
\texttt{\textbackslash hypertarget}, where \textit{slug} is the shard identifier with its
first hyphen replaced by a colon --- \texttt{def-co-top} becomes \texttt{def:co-top}. That
is the same identifier-to-label transform the provenance gate uses for results, so a lemma
section may write \texttt{\textbackslash ref\{def:co-top\}} and the knowledge-graph page
may anchor at the same target. Inside the generated statements, a crosslink is a live link
exactly when its target is in this document --- a definition rendered in this section, or a
registry result reproduced in a later one --- and otherwise is printed as a bare identifier,
because a link into material this document does not contain would be a promise it cannot
keep. Each definition also names its canonical shard path, so the reader can go from the
page to the source of truth in one step, and from there to
\texttt{definitions/INDEX.md} and \texttt{argument/DAG.md}.

% GENERATED by scripts/gen-report-defs.py — DO NOT HAND-EDIT.
% Single source of truth: definitions/<id>.md (CLAUDE.md L2).  Re-render with
%   python3 scripts/gen-report-defs.py
% The gate `python3 scripts/gen-report-defs.py --check` fails when this file is stale.

% --- rendering macros for the generated definitions (\providecommand: report/main.tex may
% --- override them; the visual language mirrors \contractquote there) -------------------
\ifdefined\defsourcequote\else
  \newenvironment{defsourcequote}{%
    \par\addvspace{0.5\baselineskip}\footnotesize
    \noindent\hrulefill\par\nobreak\vspace{0.35ex}%
    \ttfamily\frenchspacing\raggedright\sloppy}%
   {\par\vspace{0.35ex}\noindent\hrulefill\par\addvspace{0.5\baselineskip}}
\fi
\providecommand{\defsourceline}[1]{\noindent\detokenize{#1}\par}
\providecommand{\defsourcegap}{\par\addvspace{0.4\baselineskip}}
\ifdefined\defnote\else
  \newenvironment{defnote}{%
    \par\addvspace{0.35\baselineskip}\footnotesize\raggedright\sloppy
    \noindent\ignorespaces}%
   {\par\addvspace{0.5\baselineskip}}
\fi

% The classical picture: the signed and stochastic vocabulary
% GENERATED by scripts/gen-report-defs.py — DO NOT HAND-EDIT.
% Single source of truth: definitions/<id>.md (CLAUDE.md L2).  Re-render with
%   python3 scripts/gen-report-defs.py
% The gate `python3 scripts/gen-report-defs.py --check` fails when this file is stale.

% Layer file: The classical picture: the signed and stochastic vocabulary
\subsection*{The classical picture: the signed and stochastic vocabulary}

%% ---- def-negative-mass -------------------------------------------------
\hypertarget{def:negative-mass}{}%
\begin{definition}[negative mass\normalfont{} (\texttt{def-\allowbreak{}negative-\allowbreak{}mass})]\label{def:negative-mass}
\emph{Aliases:} negativity; row polytope; row polytope K\par\smallskip
\textbf{Statement.} For an \hyperref[def:signed-idempotent]{exact signed idempotent} $P\in\mathbb R^{n\times n}$ (or any
signed matrix), the \emph{negative mass} (equivalently \emph{negativity}) is
\[\delta(P):=\max_{1\le i\le n}\ \sum_{j=1}^{n}\max\{-P_{ij},\,0\},\]
the largest total negative mass in any single row. The \emph{row polytope} is
\[K:=\operatorname{conv}\{p_1,\dots,p_n\},\]
the convex hull of the rows $p_i\in\mathbb R^n$. A \hyperref[def:stochastic]{stochastic idempotent} is exactly a
signed idempotent with $\delta(P)=0$.

\textbf{Notes / provenance.} $\delta(P)$ quoted from \texttt{classical-\allowbreak{}portfolio/\allowbreak{}report/\allowbreak{}kernel-\allowbreak{}conjecture.\allowbreak{}tex}
\S{}Setting (Definition: \emph{Exact signed idempotent}); the row polytope $K$ adopted from
\texttt{..\allowbreak{}/\allowbreak{}almost-\allowbreak{}idempotent-\allowbreak{}positive-\allowbreak{}maps/\allowbreak{}definitions/\allowbreak{}def-\allowbreak{}stochastic.\allowbreak{}md}. This scalar is the single knob that
controls the geometry: the three derived scales $\tau=\sqrt\delta$, $\rho=4\tau$, $\kappa=\tau/4$ live in
\hyperref[def:visible-set]{\texttt{def-\allowbreak{}visible-\allowbreak{}set}}. The $\delta$-positivity of \hyperref[def:near-positive-projection]{\texttt{def-\allowbreak{}near-\allowbreak{}positive-\allowbreak{}projection}} is the same quantity: a
unital idempotent $R$ is $\delta$-positive iff $\delta(R)\le\delta$.
\end{definition}
\begin{defnote}
\textbf{Kind.} original (introduced by this project); \texttt{status: locked}.
\par
\textbf{Provenance.} \texttt{source: internal}; locus \texttt{quoted from classical-\allowbreak{}portfolio/\allowbreak{}report/\allowbreak{}kernel-\allowbreak{}conjecture.\allowbreak{}tex \S{}Setting (def:\allowbreak{}esi,\allowbreak{} \textnormal{$\delta$}(P));\allowbreak{} row polytope from ..\allowbreak{}/\allowbreak{}almost-\allowbreak{}idempotent-\allowbreak{}positive-\allowbreak{}maps/\allowbreak{}definitions/\allowbreak{}def-\allowbreak{}stochastic.\allowbreak{}md}; no source hash (not a literature transcription); record: adapted from kernel-\allowbreak{}conjecture.\allowbreak{}tex \S{}Setting (negativity \textnormal{$\delta$}(P)) and ..\allowbreak{}/\allowbreak{}almost-\allowbreak{}idempotent-\allowbreak{}positive-\allowbreak{}maps definitions/\allowbreak{}def-\allowbreak{}stochastic.\allowbreak{}md (row polytope K).
\par
\textbf{Used in this report by} \hyperref[lem:classical-equiv]{\texttt{lem-\allowbreak{}classical-\allowbreak{}equiv}} (164 registry results import it in total; see \texttt{argument/INDEX.md}).
\par
\textbf{Canonical shard.} \texttt{definitions/\allowbreak{}def-\allowbreak{}negative-\allowbreak{}mass.\allowbreak{}md}; twins: \texttt{definitions/INDEX.md}, \texttt{argument/DAG.md}.
\end{defnote}

%% ---- def-signed-idempotent ---------------------------------------------
\hypertarget{def:signed-idempotent}{}%
\begin{definition}[exact signed idempotent\normalfont{} (\texttt{def-\allowbreak{}signed-\allowbreak{}idempotent})]\label{def:signed-idempotent}
\emph{Aliases:} signed idempotent; signed affine retraction; signed stochastic retraction\par\smallskip
\textbf{Statement (exact signed idempotent).} Let $n\in\mathbb N$. A matrix $P\in\mathbb R^{n\times n}$ is an
\emph{exact signed idempotent} if
\[P\mathbf 1=\mathbf 1\qquad\text{and}\qquad P^2=P,\]
where $\mathbf 1=(1,\dots,1)^{\mathsf T}$. Equivalently, $P$ is a linear self-map of $\mathbb R^n$ fixing
the all-ones vector with $P^2=P$ \emph{exactly}, whose rows are signed measures of total mass $1$
($\sum_j P_{ij}=1$ for each $i$). A row-stochastic idempotent --- a \hyperref[def:stochastic]{stochastic idempotent}
--- is exactly an exact signed idempotent whose \hyperref[def:negative-mass]{negative mass} $\delta(P)=0$.

\emph{Row geometry.} Writing $p_i\in\mathbb R^n$ for the $i$-th row with the $\ell^1$ metric, each row
satisfies $\sum_j p_{ij}=1$ and $\sum_j|p_{ij}|\le1+2\delta$; hence all rows lie in a common affine
hyperplane and pairwise $\ell^1$ distances are at most $2+4\delta$ (with $\delta=\delta(P)$).

\textbf{Notes / provenance.} Quoted from \texttt{classical-\allowbreak{}portfolio/\allowbreak{}report/\allowbreak{}kernel-\allowbreak{}conjecture.\allowbreak{}tex} \S{}Setting
(Definition: \emph{Exact signed idempotent}, and the row-geometry remark). The exact-idempotence form is the
working object: $P^2=P$ holds exactly while non-positivity is quarantined into the scalar
\hyperref[def:negative-mass]{negative mass} $\delta(P)$. This is the matrix/geometry framing of
\hyperref[def:near-positive-projection]{\texttt{def-\allowbreak{}near-\allowbreak{}positive-\allowbreak{}projection}} (map/positivity framing). Its exposed row vertices, visible set, height
and invisible mass are \hyperref[def:exposed]{\texttt{def-\allowbreak{}exposed}}, \hyperref[def:visible-set]{\texttt{def-\allowbreak{}visible-\allowbreak{}set}}, \hyperref[def:height]{\texttt{def-\allowbreak{}height}}, \hyperref[def:invisible-mass]{\texttt{def-\allowbreak{}invisible-\allowbreak{}mass}}.
\end{definition}
\begin{defnote}
\textbf{Kind.} original (introduced by this project); \texttt{status: locked}.
\par
\textbf{Provenance.} \texttt{source: internal}; locus \texttt{quoted from classical-\allowbreak{}portfolio/\allowbreak{}report/\allowbreak{}kernel-\allowbreak{}conjecture.\allowbreak{}tex \S{}Setting (Definition,\allowbreak{} exact signed idempotent)}; no source hash (not a literature transcription); record: adapted from ..\allowbreak{}/\allowbreak{}almost-\allowbreak{}idempotent-\allowbreak{}positive-\allowbreak{}maps/\allowbreak{}agent-\allowbreak{}A/\allowbreak{}explorations/\allowbreak{}classical-\allowbreak{}portfolio/\allowbreak{}report/\allowbreak{}kernel-\allowbreak{}conjecture.\allowbreak{}tex \S{}Setting (def:\allowbreak{}esi);\allowbreak{} exact-\allowbreak{}idempotence form adopted for this repo.
\par
\textbf{Used in this report by} \hyperref[lem:classical-equiv]{\texttt{lem-\allowbreak{}classical-\allowbreak{}equiv}} (186 registry results import it in total; see \texttt{argument/INDEX.md}).
\par
\textbf{Canonical shard.} \texttt{definitions/\allowbreak{}def-\allowbreak{}signed-\allowbreak{}idempotent.\allowbreak{}md}; twins: \texttt{definitions/INDEX.md}, \texttt{argument/DAG.md}.
\end{defnote}

%% ---- def-stochastic ----------------------------------------------------
\hypertarget{def:stochastic}{}%
\begin{definition}[stochastic idempotent\normalfont{} (\texttt{def-\allowbreak{}stochastic})]\label{def:stochastic}
\emph{Aliases:} row-stochastic idempotent; stochastic retraction; unital positive map on $\ell\infty$\_n\par\smallskip
\textbf{Statement (commutative vocabulary).} A \emph{unital positive map} of $\ell^\infty_n=\mathbb R^n$ is a
\emph{row-stochastic matrix} $Q$ ($Q\ge0$ entrywise, $Q\mathbf 1=\mathbf 1$) --- an affine self-map of the
probability simplex $\Delta_n$. A \emph{stochastic idempotent} $E$ is a row-stochastic $E$ with $E^2=E$: an
affine retraction of $\Delta_n$ onto a sub-polytope.

Dropping positivity but keeping \emph{exact} idempotence gives the \emph{signed} generalization --- a
\hyperref[def:signed-idempotent]{signed affine retraction} $P$ with $P\mathbf 1=\mathbf 1$ and $P^2=P$ exactly,
whose rows are signed measures of total mass $1$. Its failure of positivity is quantified by the
\hyperref[def:negative-mass]{negative mass} $\delta(P)$; a stochastic idempotent is exactly a signed idempotent
with $\delta(P)=0$.

\textbf{Notes / provenance.} Adopted from \texttt{..\allowbreak{}/\allowbreak{}almost-\allowbreak{}idempotent-\allowbreak{}positive-\allowbreak{}maps/\allowbreak{}definitions/\allowbreak{}def-\allowbreak{}stochastic.\allowbreak{}md}
(B's classical formulation, report sec:classical \texttt{def:\allowbreak{}stochastic}); re-tagged \texttt{consensus} /
\texttt{source:\allowbreak{} internal} for this repo. The signed form is the working object because $P^2=P$ holds \emph{exactly}
(non-positivity quarantined in $\delta$, see \hyperref[def:negative-mass]{\texttt{def-\allowbreak{}negative-\allowbreak{}mass}}). This is the commutative
specialization of \hyperref[def:near-positive-projection]{\texttt{def-\allowbreak{}near-\allowbreak{}positive-\allowbreak{}projection}}; the approximate-idempotence sibling is
\hyperref[def:almost-idempotent]{\texttt{def-\allowbreak{}almost-\allowbreak{}idempotent}} (a genuinely positive $Q$ with $\lVert Q^2-Q\rVert$ small). Exposed row
vertices and the visible set are \hyperref[def:exposed]{\texttt{def-\allowbreak{}exposed}} / \hyperref[def:visible-set]{\texttt{def-\allowbreak{}visible-\allowbreak{}set}}.
\end{definition}
\begin{defnote}
\textbf{Kind.} consensus (project-internal, agreed; no single external source); \texttt{status: locked}.
\par
\textbf{Provenance.} \texttt{source: internal}; locus \texttt{adopted from ..\allowbreak{}/\allowbreak{}almost-\allowbreak{}idempotent-\allowbreak{}positive-\allowbreak{}maps/\allowbreak{}definitions/\allowbreak{}def-\allowbreak{}stochastic.\allowbreak{}md}; no source hash (not a literature transcription); record: adopted from ..\allowbreak{}/\allowbreak{}almost-\allowbreak{}idempotent-\allowbreak{}positive-\allowbreak{}maps definitions/\allowbreak{}def-\allowbreak{}stochastic.\allowbreak{}md (B's classical formulation,\allowbreak{} report sec:\allowbreak{}classical def:\allowbreak{}stochastic).
\par
\textbf{Used in this report by} \hyperref[lem:classical-equiv]{\texttt{lem-\allowbreak{}classical-\allowbreak{}equiv}}, \hyperref[lem:prh]{\texttt{lem-\allowbreak{}prh}}, \hyperref[lem:routef-prh-finish]{\texttt{lem-\allowbreak{}routef-\allowbreak{}prh-\allowbreak{}finish}} (17 registry results import it in total; see \texttt{argument/INDEX.md}).
\par
\textbf{Canonical shard.} \texttt{definitions/\allowbreak{}def-\allowbreak{}stochastic.\allowbreak{}md}; twins: \texttt{definitions/INDEX.md}, \texttt{argument/DAG.md}.
\end{defnote}

%% ---- def-almost-idempotent ---------------------------------------------
\hypertarget{def:almost-idempotent}{}%
\begin{definition}[almost idempotent\normalfont{} (\texttt{def-\allowbreak{}almost-\allowbreak{}idempotent})]\label{def:almost-idempotent}
\emph{Aliases:} $\eta$-idempotent; eta-idempotent; almost-idempotent stochastic matrix\par\smallskip
\textbf{Statement (commutative / stochastic specialization).} A \hyperref[def:stochastic]{row-stochastic matrix} $Q$
on $\ell^\infty_n$ is \emph{almost idempotent} (with defect $\eta$) if
\[\lVert Q^2-Q\rVert_{\infty\to\infty}\le\eta,\qquad \eta\in[0,\tfrac14),\]
where $\lVert\cdot\rVert_{\infty\to\infty}$ is the operator norm induced by $\ell^\infty_n$ --- i.e. the
maximum over rows of the $\ell^1$-norm of the corresponding row of $Q^2-Q$. The range $\eta<\tfrac14$ is
what makes the binomial series producing the exact idempotent converge.

\textbf{Notes / provenance.} Adopted from \texttt{..\allowbreak{}/\allowbreak{}almost-\allowbreak{}idempotent-\allowbreak{}positive-\allowbreak{}maps/\allowbreak{}definitions/\allowbreak{}def-\allowbreak{}almost-\allowbreak{}idempotent.\allowbreak{}md}
(the operator-norm --- \emph{not} cb-norm --- formulation, report rem:cb-norm), specialized to the
$\lVert\cdot\rVert_{\infty\to\infty}$ map norm on stochastic matrices as in
\texttt{classical-\allowbreak{}portfolio/\allowbreak{}report/\allowbreak{}kernel-\allowbreak{}conjecture.\allowbreak{}tex} Theorem (chain)(c). Measuring non-idempotence in the
map operator norm (rather than the cb-norm) is why the bridge exponent is $\sqrt\eta$ rather than $\eta$.
The classical stability question (\texttt{op-\allowbreak{}classical}) asks whether every such $Q$ lies within $C\sqrt\eta$
(max-row-$\ell^1$) of an exactly idempotent row-stochastic matrix; its parent object is
\hyperref[def:near-positive-projection]{\texttt{def-\allowbreak{}near-\allowbreak{}positive-\allowbreak{}projection}}. Distinguish $\eta$ (this defect) from the
\hyperref[def:negative-mass]{negative mass} $\delta$ of the associated \hyperref[def:signed-idempotent]{\texttt{def-\allowbreak{}signed-\allowbreak{}idempotent}}, with
$\delta=O(\eta)$.
\end{definition}
\begin{defnote}
\textbf{Kind.} consensus (project-internal, agreed; no single external source); \texttt{status: locked}.
\par
\textbf{Provenance.} \texttt{source: internal}; locus \texttt{adopted from ..\allowbreak{}/\allowbreak{}almost-\allowbreak{}idempotent-\allowbreak{}positive-\allowbreak{}maps/\allowbreak{}definitions/\allowbreak{}def-\allowbreak{}almost-\allowbreak{}idempotent.\allowbreak{}md (stochastic specialization)}; no source hash (not a literature transcription); record: adopted from ..\allowbreak{}/\allowbreak{}almost-\allowbreak{}idempotent-\allowbreak{}positive-\allowbreak{}maps definitions/\allowbreak{}def-\allowbreak{}almost-\allowbreak{}idempotent.\allowbreak{}md (operator-\allowbreak{}norm formulation,\allowbreak{} report rem:\allowbreak{}cb-\allowbreak{}norm);\allowbreak{} \textnormal{$\infty\to\infty$} specialization per kernel-\allowbreak{}conjecture.\allowbreak{}tex thm:\allowbreak{}chain(c).
\par
\textbf{Used in this report by} \hyperref[lem:classical-equiv]{\texttt{lem-\allowbreak{}classical-\allowbreak{}equiv}} (4 registry results import it in total; see \texttt{argument/INDEX.md}).
\par
\textbf{Canonical shard.} \texttt{definitions/\allowbreak{}def-\allowbreak{}almost-\allowbreak{}idempotent.\allowbreak{}md}; twins: \texttt{definitions/INDEX.md}, \texttt{argument/DAG.md}.
\end{defnote}

%% ---- def-exposed -------------------------------------------------------
\hypertarget{def:exposed}{}%
\begin{definition}[exposed vertex\normalfont{} (\texttt{def-\allowbreak{}exposed})]\label{def:exposed}
\emph{Aliases:} ($\rho$,$\kappa$)-exposed; (rho,kappa)-exposed; exposedness margin; admissible exposer; hidden row vertex; well-exposed vertex\par\smallskip
\textbf{Statement.} Fix an \hyperref[def:signed-idempotent]{exact signed idempotent} $P\in\mathbb R^{n\times n}$ and
write $p_i\in\mathbb R^n$ for its $i$-th row, with the $\ell^1$ metric $\lVert x-y\rVert_1=\sum_j|x_j-y_j|$.

A row $p_v$ is a \emph{row vertex} of $P$ if $p_v\notin\operatorname{conv}\{\,p_j:p_j\neq p_v\text{ as vectors of }\mathbb R^n\,\}$
(geometrically coincident duplicate rows count as a single point; a repeated point is still a vertex if it
lies outside the hull of the \emph{other} points).

An \emph{admissible exposer} for a row vertex $v$ is an affine function $h:\mathbb R^n\to\mathbb R$ with
\[h(p_v)=0\qquad\text{and}\qquad 0\le h(p_j)\le 1\quad\text{for every row }p_j.\]
The \emph{exposedness margin} of $v$ is
\[t^*(v):=\sup_{h\text{ admissible}}\ \min\{\,h(p_j):\lVert p_j-p_v\rVert_1\ge\rho\,\},\]
with the convention $t^*(v)=+\infty$ when no row is at $\ell^1$-distance $\ge\rho$ from $p_v$. The vertex
$v$ is \emph{$(\rho,\kappa)$-exposed} if $t^*(v)\ge\kappa$, and \emph{hidden} otherwise. The scales $\rho=4\tau$,
$\kappa=\tau/4$ (with $\tau=\sqrt\delta$) and the resulting \emph{visible set} $\mathcal W(P)$ of exposed
vertices are recorded in \hyperref[def:visible-set]{\texttt{def-\allowbreak{}visible-\allowbreak{}set}}.

Informally: an exposed vertex can be linearly separated, with margin $\ge\kappa$, from every row that is
genuinely far ($\ge\rho$) from it; rows inside the $\rho$-ball are exempt.

\textbf{Notes / provenance.} Adapted from \texttt{classical-\allowbreak{}portfolio/\allowbreak{}report/\allowbreak{}kernel-\allowbreak{}conjecture.\allowbreak{}tex} \S{}Setting
(def:vertex, def:exposed) and \texttt{..\allowbreak{}/\allowbreak{}almost-\allowbreak{}idempotent-\allowbreak{}positive-\allowbreak{}maps/\allowbreak{}definitions/\allowbreak{}def-\allowbreak{}exposed.\allowbreak{}md}. The latter
states the equivalent \emph{exposedness modulus} $e_v(\rho)=\sup_h\min_{\lVert p_i-v\rVert_1\ge\rho}h(p_i)$ and
calls a vertex \emph{well-exposed} at $\sqrt{}$-scale if $e_v(\rho)\ge c\sqrt\delta$ for some
$\rho=O(\sqrt\delta)$, where $\delta$ is the \hyperref[def:negative-mass]{negative mass}. Central to the proved
classical special cases (well-exposed $\Rightarrow$ simplex) and to the open global exposed-hull lemma
(\texttt{op-\allowbreak{}exposed-\allowbreak{}hull}). A pointwise exposed-or-redundant dichotomy is provably \emph{insufficient} (dense regular
polygons) --- the gap must be stated globally (see \hyperref[def:height]{\texttt{def-\allowbreak{}height}}).
\end{definition}
\begin{defnote}
\textbf{Kind.} consensus (project-internal, agreed; no single external source); \texttt{status: locked}.
\par
\textbf{Provenance.} \texttt{source: internal}; locus \texttt{adapted from classical-\allowbreak{}portfolio/\allowbreak{}report/\allowbreak{}kernel-\allowbreak{}conjecture.\allowbreak{}tex \S{}Setting (def:\allowbreak{}vertex,\allowbreak{} def:\allowbreak{}exposed);\allowbreak{} ..\allowbreak{}/\allowbreak{}almost-\allowbreak{}idempotent-\allowbreak{}positive-\allowbreak{}maps/\allowbreak{}definitions/\allowbreak{}def-\allowbreak{}exposed.\allowbreak{}md}; no source hash (not a literature transcription); record: adopted from ..\allowbreak{}/\allowbreak{}almost-\allowbreak{}idempotent-\allowbreak{}positive-\allowbreak{}maps definitions/\allowbreak{}def-\allowbreak{}exposed.\allowbreak{}md (B's exposed-\allowbreak{}circuit machinery,\allowbreak{} report sec:\allowbreak{}classical);\allowbreak{} reconciled with kernel-\allowbreak{}conjecture.\allowbreak{}tex def:\allowbreak{}exposed.
\par
\textbf{Used by} 129 registry results, none of them reproduced in this report (exploration track; see \texttt{argument/INDEX.md} and \texttt{report/UNWIRED.md}).
\par
\textbf{Canonical shard.} \texttt{definitions/\allowbreak{}def-\allowbreak{}exposed.\allowbreak{}md}; twins: \texttt{definitions/INDEX.md}, \texttt{argument/DAG.md}.
\end{defnote}

%% ---- def-visible-set ---------------------------------------------------
\hypertarget{def:visible-set}{}%
\begin{definition}[visible set\normalfont{} (\texttt{def-\allowbreak{}visible-\allowbreak{}set})]\label{def:visible-set}
\emph{Aliases:} visible set W; W(P); ($\rho$,$\kappa$)-exposed set; three scales\par\smallskip
\textbf{Statement.} Given an \hyperref[def:signed-idempotent]{exact signed idempotent} $P$ with
$\delta=\delta(P)$ (the \hyperref[def:negative-mass]{negative mass}), the \emph{three scales} are
\[\tau:=\sqrt\delta,\qquad \rho:=4\tau,\qquad \kappa:=\tau/4.\]
The \emph{visible set} is
\[\mathcal W=\mathcal W(P):=\{\,v:\ p_v\text{ is a }(\rho,\kappa)\text{-exposed row vertex of }P\,\},\]
the set of row vertices that are \hyperref[def:exposed]{$(\rho,\kappa)$-exposed} at these scales; all other row
vertices are \emph{hidden}. At $\delta=0$ the visible vertices are exactly the distinct rows of the recurrent
(equal-input) blocks of a stochastic idempotent.

\textbf{Notes / provenance.} Adapted from \texttt{classical-\allowbreak{}portfolio/\allowbreak{}report/\allowbreak{}kernel-\allowbreak{}conjecture.\allowbreak{}tex} \S{}Setting
(Definition: \emph{The three scales}; Definition: \emph{$(\rho,\kappa)$-exposedness and the visible set $\mathcal W$}).
The scales couple exposedness (\hyperref[def:exposed]{\texttt{def-\allowbreak{}exposed}}) to the negative mass (\hyperref[def:negative-mass]{\texttt{def-\allowbreak{}negative-\allowbreak{}mass}}). The convex
hull $C_{\mathcal W}:=\operatorname{conv}\{p_w:w\in\mathcal W\}$ is consumed by \hyperref[def:height]{\texttt{def-\allowbreak{}height}} (max
$\ell^1$-distance of a row from $C_{\mathcal W}$) and \hyperref[def:invisible-mass]{\texttt{def-\allowbreak{}invisible-\allowbreak{}mass}} (positive mass a row places
outside $C_{\mathcal W}$).
\end{definition}
\begin{defnote}
\textbf{Kind.} original (introduced by this project); \texttt{status: locked}.
\par
\textbf{Provenance.} \texttt{source: internal}; locus \texttt{quoted from classical-\allowbreak{}portfolio/\allowbreak{}report/\allowbreak{}kernel-\allowbreak{}conjecture.\allowbreak{}tex \S{}Setting (def:\allowbreak{}scales,\allowbreak{} def:\allowbreak{}exposed)}; no source hash (not a literature transcription); record: adapted from kernel-\allowbreak{}conjecture.\allowbreak{}tex \S{}Setting (Definition:\allowbreak{} three scales;\allowbreak{} Definition:\allowbreak{} (\textnormal{$\rho$},\textnormal{$\kappa$})-\allowbreak{}exposedness and the visible set W).
\par
\textbf{Used by} 137 registry results, none of them reproduced in this report (exploration track; see \texttt{argument/INDEX.md} and \texttt{report/UNWIRED.md}).
\par
\textbf{Canonical shard.} \texttt{definitions/\allowbreak{}def-\allowbreak{}visible-\allowbreak{}set.\allowbreak{}md}; twins: \texttt{definitions/INDEX.md}, \texttt{argument/DAG.md}.
\end{defnote}

%% ---- def-height --------------------------------------------------------
\hypertarget{def:height}{}%
\begin{definition}[height\normalfont{} (\texttt{def-\allowbreak{}height})]\label{def:height}
\emph{Aliases:} height H; H(P); hidden top vertex\par\smallskip
\textbf{Statement.} Let $C_{\mathcal W}:=\operatorname{conv}\{p_w:w\in\mathcal W\}$, where $\mathcal W=\mathcal W(P)$
is the \hyperref[def:visible-set]{\texttt{def-\allowbreak{}visible-\allowbreak{}set}}, and let $\operatorname{dist}_1(x,S)=\inf_{s\in\operatorname{conv}S}\lVert x-s\rVert_1$.
If $\mathcal W\neq\emptyset$, the \emph{height} of a row $p_i$ is $\operatorname{dist}_1(p_i,C_{\mathcal W})$, and
the \emph{height of $P$} is
\[H=H(P):=\max_{1\le i\le n}\ \operatorname{dist}_1\!\big(p_i,\ \operatorname{conv}\{p_w:w\in\mathcal W\}\big).\]
Because $\operatorname{dist}_1(\cdot,C_{\mathcal W})$ is convex and every row is a convex combination of the
geometrically distinct row vertices, the maximum is attained at a row vertex; if $H>0$, any maximizing
vertex is necessarily \hyperref[def:exposed]{hidden}. Such a vertex is called a \emph{hidden top vertex}.

\textbf{Notes / provenance.} Quoted from \texttt{classical-\allowbreak{}portfolio/\allowbreak{}report/\allowbreak{}kernel-\allowbreak{}conjecture.\allowbreak{}tex} \S{}Setting
(Definition: \emph{Height}). $H$ measures how far the rows spill outside the hull of the visible
(\hyperref[def:visible-set]{\texttt{def-\allowbreak{}visible-\allowbreak{}set}}, \hyperref[def:exposed]{\texttt{def-\allowbreak{}exposed}}) vertices; the linear hull bound / HLC target is
$H(P)\le C\sqrt\delta$ in the \hyperref[def:negative-mass]{negative mass} $\delta$. The invisible mass
\hyperref[def:invisible-mass]{\texttt{def-\allowbreak{}invisible-\allowbreak{}mass}} of a hidden top vertex is the hypothesis quantity that governs whether $H$ can be
large.
\end{definition}
\begin{defnote}
\textbf{Kind.} original (introduced by this project); \texttt{status: locked}.
\par
\textbf{Provenance.} \texttt{source: internal}; locus \texttt{quoted from classical-\allowbreak{}portfolio/\allowbreak{}report/\allowbreak{}kernel-\allowbreak{}conjecture.\allowbreak{}tex \S{}Setting (Definition:\allowbreak{} Height)}; no source hash (not a literature transcription); record: adapted from kernel-\allowbreak{}conjecture.\allowbreak{}tex \S{}Setting (Definition:\allowbreak{} Height,\allowbreak{} def:\allowbreak{}height).
\par
\textbf{Used by} 112 registry results, none of them reproduced in this report (exploration track; see \texttt{argument/INDEX.md} and \texttt{report/UNWIRED.md}).
\par
\textbf{Canonical shard.} \texttt{definitions/\allowbreak{}def-\allowbreak{}height.\allowbreak{}md}; twins: \texttt{definitions/INDEX.md}, \texttt{argument/DAG.md}.
\end{defnote}

%% ---- def-actor-hull ----------------------------------------------------
\hypertarget{def:actor-hull}{}%
\begin{definition}[always-tight hull\normalfont{} (\texttt{def-\allowbreak{}actor-\allowbreak{}hull})]\label{def:actor-hull}
\emph{Aliases:} actor hull; K\_T; K\_O; K\_T(u); K\_O(u); always-tight hulls; disjoint always-tight hulls; displacement hull\par\smallskip
\textbf{Statement.} Fix an \hyperref[def:signed-idempotent]{exact signed idempotent} $P$ and a
\hyperref[def:exposed]{hidden} geometrically distinct row vertex $u$ with $t^*(u)>0$. Let $T(u)$, $O(u)$,
$Z(u)$ be the always-tight $\rho$-far / upper-box / \hyperref[def:zero-face]{zero-face} constraint families of
the exposedness LP at $u$ (\hyperlink{dag:lem-always-tight-dual-support}{\texttt{lem-\allowbreak{}always-\allowbreak{}tight-\allowbreak{}dual-\allowbreak{}support}}). The \textbf{always-tight hulls} (or \emph{actor
hulls}) are the displacement convex bodies
\[K_T(u):=\operatorname{conv}\{p_f-p_u:f\in T(u)\},\qquad
  K_O(u):=t^*(u)\cdot\operatorname{conv}\{p_i-p_u:i\in O(u)\}.\]
They are \textbf{disjoint} when $g:=\operatorname{dist}_1(K_T(u),K_O(u))>0$ (the huddle / Branch-I
predicate), and \textbf{intersect} otherwise (Branch II, the alpha-free optimal display of
\hyperlink{dag:lem-optimal-face-conic-reduction}{\texttt{lem-\allowbreak{}optimal-\allowbreak{}face-\allowbreak{}conic-\allowbreak{}reduction}}). Boundary ownership: tangency / any common point counts as
intersecting.

\textbf{Notes / provenance.} Project-original; ``actor hull'' and ``always-tight hull(s)'' name the same
displacement bodies $K_T,K_O$ built from the actors (far / upper-box rows) the
\hyperref[def:dual-witness]{dual witness} charges. The disjoint-vs-intersect dichotomy on $K_T,K_O$ is the
root split $S1$ of the W54 huddle-charge tree. \texttt{status:\allowbreak{} draft} --- A+B sign-off pending (Rule 7).
Related: \hyperref[def:zero-face]{\texttt{def-\allowbreak{}zero-\allowbreak{}face}}, \hyperref[def:dual-witness]{\texttt{def-\allowbreak{}dual-\allowbreak{}witness}}, \hyperref[def:co-top]{\texttt{def-\allowbreak{}co-\allowbreak{}top}}.
\end{definition}
\begin{defnote}
\textbf{Kind.} original (introduced by this project); \texttt{status: draft}.
\par
\textbf{Provenance.} \texttt{source: internal}; locus \texttt{internal;\allowbreak{} first pinned in argument/\allowbreak{}lemmas/\allowbreak{}lem-\allowbreak{}downhill-\allowbreak{}cotop-\allowbreak{}conic-\allowbreak{}mass.\allowbreak{}md and lem-\allowbreak{}always-\allowbreak{}tight-\allowbreak{}dual-\allowbreak{}support.\allowbreak{}md}; no source hash (not a literature transcription); record: project-\allowbreak{}original;\allowbreak{} sign-\allowbreak{}off pending (Rule 7 lock).\allowbreak{} First pinned by lem-\allowbreak{}always-\allowbreak{}tight-\allowbreak{}dual-\allowbreak{}support (the T/\allowbreak{}O families) + lem-\allowbreak{}downhill-\allowbreak{}cotop-\allowbreak{}conic-\allowbreak{}mass (the hulls K\_T,\allowbreak{} K\_O).
\par
\textbf{Used by} 27 registry results, none of them reproduced in this report (exploration track; see \texttt{argument/INDEX.md} and \texttt{report/UNWIRED.md}).
\par
\textbf{Canonical shard.} \texttt{definitions/\allowbreak{}def-\allowbreak{}actor-\allowbreak{}hull.\allowbreak{}md}; twins: \texttt{definitions/INDEX.md}, \texttt{argument/DAG.md}.
\end{defnote}

%% ---- def-dual-witness --------------------------------------------------
\hypertarget{def:dual-witness}{}%
\begin{definition}[hiddenness dual witness\normalfont{} (\texttt{def-\allowbreak{}dual-\allowbreak{}witness})]\label{def:dual-witness}
\emph{Aliases:} dual witness; dual certificate; small-beta witness; reduced optimal witness; witness balance; (lambda, alpha, beta)\par\smallskip
\textbf{Statement.} Fix an \hyperref[def:signed-idempotent]{exact signed idempotent} $P$ and a
\hyperref[def:exposed]{hidden} row vertex $v$ ($\rho=4\tau$, $\kappa=\tau/4$, $\tau=\sqrt{\delta(P)}$), with
$F_v=\{j:\lVert p_j-p_v\rVert_1\ge\rho\}$ the (nonempty) $\rho$-far row set.

A \textbf{hiddenness dual witness} of $v$ is a triple $(\lambda,\alpha,\beta)$ of nonnegative
coefficients --- $\lambda_f\ (f\in F_v)$ with $\sum_f\lambda_f=1$ (a probability on the far rows), and
$\alpha_i,\beta_i\ge0$ over all rows with $\sum_i\beta_i=t^*(v)<\kappa$ --- satisfying the \emph{witness
balance}
\[\sum_{f\in F_v}\lambda_f\,(p_f-p_v)\;+\;\sum_i\alpha_i\,(p_i-p_v)\;=\;\sum_i\beta_i\,(p_i-p_v).\]
It is the LP dual of the exposedness program (see \hyperref[def:exposed]{\texttt{def-\allowbreak{}exposed}}); the \emph{small-beta} property is
$\sum_i\beta_i=t^*(v)<\kappa$. A witness is \textbf{reduced} when redundant centered-zero constraints are
deleted; by \hyperlink{dag:lem-always-tight-dual-support}{\texttt{lem-\allowbreak{}always-\allowbreak{}tight-\allowbreak{}dual-\allowbreak{}support}} a reduced optimal witness has
$\operatorname{supp}(\lambda)\subseteq T$, $\operatorname{supp}(\beta)\subseteq O$, and
$\operatorname{supp}(\alpha)\subseteq Z$ (the \hyperref[def:actor-hull]{always-tight} far / upper-box /
\hyperref[def:zero-face]{zero-face} families).

\textbf{Notes / provenance.} Project-original; from \hyperlink{dag:lem-hiddenness-dual-witness}{\texttt{lem-\allowbreak{}hiddenness-\allowbreak{}dual-\allowbreak{}witness}} (existence via finite
LP duality). The witness converts the scalar hiddenness fact $t^*(v)<\kappa$ into a geometric object:
a convex combination of $\rho$-far rows reproducing $p_v$ up to controlled signed slack. ``dual
certificate'' is used interchangeably in the shards. \texttt{status:\allowbreak{} draft} --- A+B sign-off pending (Rule 7).
Related: \hyperref[def:zero-face]{\texttt{def-\allowbreak{}zero-\allowbreak{}face}}, \hyperref[def:actor-hull]{\texttt{def-\allowbreak{}actor-\allowbreak{}hull}}, \hyperref[def:co-top]{\texttt{def-\allowbreak{}co-\allowbreak{}top}}.
\end{definition}
\begin{defnote}
\textbf{Kind.} original (introduced by this project); \texttt{status: draft}.
\par
\textbf{Provenance.} \texttt{source: internal}; locus \texttt{internal;\allowbreak{} first pinned in argument/\allowbreak{}lemmas/\allowbreak{}lem-\allowbreak{}hiddenness-\allowbreak{}dual-\allowbreak{}witness.\allowbreak{}md}; no source hash (not a literature transcription); record: project-\allowbreak{}original;\allowbreak{} sign-\allowbreak{}off pending (Rule 7 lock).\allowbreak{} First pinned by lem-\allowbreak{}hiddenness-\allowbreak{}dual-\allowbreak{}witness (the LP-\allowbreak{}dual object) + lem-\allowbreak{}always-\allowbreak{}tight-\allowbreak{}dual-\allowbreak{}support (support localization).
\par
\textbf{Used by} 1 registry result, none of them reproduced in this report (exploration track; see \texttt{argument/INDEX.md} and \texttt{report/UNWIRED.md}).
\par
\textbf{Canonical shard.} \texttt{definitions/\allowbreak{}def-\allowbreak{}dual-\allowbreak{}witness.\allowbreak{}md}; twins: \texttt{definitions/INDEX.md}, \texttt{argument/DAG.md}.
\end{defnote}

%% ---- def-zero-face -----------------------------------------------------
\hypertarget{def:zero-face}{}%
\begin{definition}[zero-face row\normalfont{} (\texttt{def-\allowbreak{}zero-\allowbreak{}face})]\label{def:zero-face}
\emph{Aliases:} zero face; zero-face family; always-tight zero-face family; Z(u); h-zero row; zero-face conic mass\par\smallskip
\textbf{Statement.} Fix an \hyperref[def:signed-idempotent]{exact signed idempotent} $P$ and a hidden
(\hyperref[def:exposed]{not $(\rho,\kappa)$-exposed}) geometrically distinct row vertex $u$, with the
exposedness LP at $u$ (the program defining $t^*(u)$, see \hyperref[def:exposed]{\texttt{def-\allowbreak{}exposed}}) and an \emph{optimal exposer}
$h^*$ (an admissible affine $h$ attaining $t^*(u)$).

A row $z$ is a \textbf{zero-face row} of $u$ if $h^*(p_z)=0$, i.e. $z$ lies on the zero face
$\{x:h^*(x)=0\}$ of the optimal exposer. The \textbf{always-tight zero-face family} $Z(u)$ is the set of
rows whose lower-box constraint $h\ge 0$ is tight on the \emph{whole} primal optimal face of the LP
(equivalently, $\operatorname{supp}(\alpha)\subseteq Z(u)$ for every reduced optimal hiddenness
\hyperref[def:dual-witness]{dual witness} $(\lambda,\alpha,\beta)$, by \hyperlink{dag:lem-always-tight-dual-support}{\texttt{lem-\allowbreak{}always-\allowbreak{}tight-\allowbreak{}dual-\allowbreak{}support}}).
The \textbf{zero-face conic mass} is $\sum_{z\in Z(u)}a_z$ for the reduced optimal display's zero-face
coefficients $a_z\ge 0$.

By \texttt{lem-\allowbreak{}zero-\allowbreak{}face-\allowbreak{}localization} every zero-face row is $\rho$-near $u$
($\lVert p_z-p_u\rVert_1<4\tau$, $\tau=\sqrt{\delta}$); if $u$ is within $4\tau$ of a
\hyperref[def:height]{hidden top} $v$ of height $H$ then additionally $\lVert p_z-p_v\rVert_1<8\tau$ and $z$
has depth $>H-8\tau$.

\textbf{Notes / provenance.} Project-original; the vocabulary in which the huddle-charge terminal node
is posed (kill or bound the zero-face conic mass). Distinguish two uses in the shards, kept
consistent here: the \emph{pointwise} zero-face row ($h^*(p_z)=0$ for one fixed optimal exposer) and the
\emph{always-tight family} $Z(u)$ (tight on the entire optimal face). $Z(u)$ is the always-tight version;
statements that fix one optimal exposer $h^*$ use the pointwise version, and the two agree on the
relative interior optimal exposer. \texttt{status:\allowbreak{} draft} --- A+B sign-off pending (Rule 7). Related:
\hyperref[def:actor-hull]{\texttt{def-\allowbreak{}actor-\allowbreak{}hull}} (the always-tight far/upper-box hulls $K_T,K_O$), \hyperref[def:dual-witness]{\texttt{def-\allowbreak{}dual-\allowbreak{}witness}}.
\end{definition}
\begin{defnote}
\textbf{Kind.} original (introduced by this project); \texttt{status: draft}.
\par
\textbf{Provenance.} \texttt{source: internal}; locus \texttt{internal;\allowbreak{} first pinned in argument/\allowbreak{}lemmas/\allowbreak{}lem-\allowbreak{}zero-\allowbreak{}face-\allowbreak{}localization.\allowbreak{}md and lem-\allowbreak{}always-\allowbreak{}tight-\allowbreak{}dual-\allowbreak{}support.\allowbreak{}md}; no source hash (not a literature transcription); record: project-\allowbreak{}original;\allowbreak{} sign-\allowbreak{}off pending (Rule 7 lock).\allowbreak{} First pinned by lem-\allowbreak{}zero-\allowbreak{}face-\allowbreak{}localization (rho-\allowbreak{}nearness) + lem-\allowbreak{}always-\allowbreak{}tight-\allowbreak{}dual-\allowbreak{}support (the alpha-\allowbreak{}carrier family Z).
\par
\textbf{Used by} 2 registry results, none of them reproduced in this report (exploration track; see \texttt{argument/INDEX.md} and \texttt{report/UNWIRED.md}).
\par
\textbf{Canonical shard.} \texttt{definitions/\allowbreak{}def-\allowbreak{}zero-\allowbreak{}face.\allowbreak{}md}; twins: \texttt{definitions/INDEX.md}, \texttt{argument/DAG.md}.
\end{defnote}

%% ---- def-co-top --------------------------------------------------------
\hypertarget{def:co-top}{}%
\begin{definition}[co-top\normalfont{} (\texttt{def-\allowbreak{}co-\allowbreak{}top})]\label{def:co-top}
\emph{Aliases:} co-top row; co-top vertex; co-top web; co-top far web; co-top zero-face mass; starved set\par\smallskip
\textbf{Statement.} Fix an \hyperref[def:signed-idempotent]{exact signed idempotent} $P$ with nonempty
\hyperref[def:visible-set]{visible set} $\mathcal W$ and a \hyperref[def:height]{hidden top vertex} $v$ of height
$H$; write $d_j=\operatorname{dist}_1(p_j,\operatorname{conv}\{p_w:w\in\mathcal W\})$,
$\tau=\sqrt{\delta(P)}$.

A row $f$ is \textbf{co-top} (relative to $v$, at slack $c$) if it sits nearly as deep as the top:
\[d_f\;>\;H-c\,\tau\qquad(c>0\ \text{a small universal constant, typically }c\in\{4,8\}).\]
The \textbf{co-top far web} of $v$ is the set of rows that are simultaneously co-top \textbf{and} $\rho$-far
from $v$ ($\lVert p_f-p_v\rVert_1\ge4\tau$) --- the \textbf{starved set}
$\{j:\lVert p_j-p_v\rVert_1\ge4\tau,\ d_j>H-8\tau\}$. By \texttt{lem-\allowbreak{}cotop-\allowbreak{}witness-\allowbreak{}pinning} more than
$13/16$ of a hidden top's \hyperref[def:dual-witness]{dual-witness} $\lambda$-mass sits in this starved set
(at $c=4$, $\delta\le1/4$), and by \texttt{lem-\allowbreak{}downhill-\allowbreak{}cotop-\allowbreak{}conic-\allowbreak{}mass} (disjoint
\hyperref[def:actor-hull]{always-tight hulls}) a definite share of the \hyperref[def:zero-face]{zero-face} conic mass
lands on nonclone, $\rho$-near, co-top rows.

\textbf{Notes / provenance.} Project-original; ``co-top'' names the deep-far actors that a hidden top's dual
witness is forced to charge yet its positive coefficient mass need not reach (the dual-required /
primal-starved tension at the heart of the huddle charge and \hyperlink{dag:conj-cotop-web-coupling}{\texttt{conj-\allowbreak{}cotop-\allowbreak{}web-\allowbreak{}coupling}}). `status:
draft` --- A+B sign-off pending (Rule 7). Related: \hyperref[def:dual-witness]{\texttt{def-\allowbreak{}dual-\allowbreak{}witness}}, \hyperref[def:actor-hull]{\texttt{def-\allowbreak{}actor-\allowbreak{}hull}},
\hyperref[def:near-cluster]{\texttt{def-\allowbreak{}near-\allowbreak{}cluster}}.
\end{definition}
\begin{defnote}
\textbf{Kind.} original (introduced by this project); \texttt{status: draft}.
\par
\textbf{Provenance.} \texttt{source: internal}; locus \texttt{internal;\allowbreak{} first pinned in argument/\allowbreak{}lemmas/\allowbreak{}lem-\allowbreak{}cotop-\allowbreak{}witness-\allowbreak{}pinning.\allowbreak{}md and lem-\allowbreak{}downhill-\allowbreak{}cotop-\allowbreak{}conic-\allowbreak{}mass.\allowbreak{}md}; no source hash (not a literature transcription); record: project-\allowbreak{}original;\allowbreak{} sign-\allowbreak{}off pending (Rule 7 lock).\allowbreak{} First pinned by lem-\allowbreak{}cotop-\allowbreak{}witness-\allowbreak{}pinning + lem-\allowbreak{}downhill-\allowbreak{}cotop-\allowbreak{}conic-\allowbreak{}mass.
\par
\textbf{Used by} 40 registry results, none of them reproduced in this report (exploration track; see \texttt{argument/INDEX.md} and \texttt{report/UNWIRED.md}).
\par
\textbf{Canonical shard.} \texttt{definitions/\allowbreak{}def-\allowbreak{}co-\allowbreak{}top.\allowbreak{}md}; twins: \texttt{definitions/INDEX.md}, \texttt{argument/DAG.md}.
\end{defnote}

%% ---- def-invisible-mass ------------------------------------------------
\hypertarget{def:invisible-mass}{}%
\begin{definition}[invisible mass\normalfont{} (\texttt{def-\allowbreak{}invisible-\allowbreak{}mass})]\label{def:invisible-mass}
\emph{Aliases:} invisible positive mass; $\widetilde\sigma$\_v; invisible mass $\widetilde\sigma$\_v\par\smallskip
\textbf{Statement.} Let $C_{\mathcal W}:=\operatorname{conv}\{p_w:w\in\mathcal W\}$ (\hyperref[def:visible-set]{\texttt{def-\allowbreak{}visible-\allowbreak{}set}}) and let
$p_v$ be a row vertex. The \emph{invisible (positive) mass} of $v$ is
\[\widetilde\sigma_v:=\sum_{j\,:\,\operatorname{dist}_1(p_j,\,C_{\mathcal W})>0}\ \max\{P_{vj},\,0\},\]
the total positive coefficient mass that row $v$ places on rows lying strictly outside $C_{\mathcal W}$ ---
\emph{including} the index $j=v$ itself when $p_v\notin C_{\mathcal W}$ and $P_{vv}>0$.

\emph{Halo robustness.} Replacing the condition $\operatorname{dist}_1(p_j,C_{\mathcal W})>0$ by
$\operatorname{dist}_1(p_j,C_{\mathcal W})>\varepsilon$ changes $\widetilde\sigma_v$ only by mass sitting
in the $\varepsilon$-halo of $C_{\mathcal W}$, which is absorbable as an $\varepsilon$-loss in the
\hyperref[def:height]{height}; the $\varepsilon=0$ form is used throughout.

\textbf{Notes / provenance.} Quoted from \texttt{classical-\allowbreak{}portfolio/\allowbreak{}report/\allowbreak{}kernel-\allowbreak{}conjecture.\allowbreak{}tex} \S{}Setting
(Definition: \emph{Invisible mass $\widetilde\sigma_v$}) with the \emph{Halo robustness} remark. $\widetilde\sigma_v$
is the hypothesis quantity of the Kernel Conjecture: a hidden top vertex (\hyperref[def:height]{\texttt{def-\allowbreak{}height}}) with
$\widetilde\sigma_v>\tau=\sqrt\delta$ is the missing branch, while the $\widetilde\sigma_v\le\tau$ branch is
the proved height cap. Uses \hyperref[def:visible-set]{\texttt{def-\allowbreak{}visible-\allowbreak{}set}}, \hyperref[def:negative-mass]{\texttt{def-\allowbreak{}negative-\allowbreak{}mass}}.
\end{definition}
\begin{defnote}
\textbf{Kind.} original (introduced by this project); \texttt{status: locked}.
\par
\textbf{Provenance.} \texttt{source: internal}; locus \texttt{quoted from classical-\allowbreak{}portfolio/\allowbreak{}report/\allowbreak{}kernel-\allowbreak{}conjecture.\allowbreak{}tex \S{}Setting (Definition:\allowbreak{} Invisible mass \textnormal{$\widetilde\sigma$}\_v;\allowbreak{} Remark:\allowbreak{} Halo robustness)}; no source hash (not a literature transcription); record: adapted from kernel-\allowbreak{}conjecture.\allowbreak{}tex \S{}Setting (Definition:\allowbreak{} Invisible mass \textnormal{$\widetilde\sigma$}\_v,\allowbreak{} def:\allowbreak{}sigt).
\par
\textbf{Used by} 35 registry results, none of them reproduced in this report (exploration track; see \texttt{argument/INDEX.md} and \texttt{report/UNWIRED.md}).
\par
\textbf{Canonical shard.} \texttt{definitions/\allowbreak{}def-\allowbreak{}invisible-\allowbreak{}mass.\allowbreak{}md}; twins: \texttt{definitions/INDEX.md}, \texttt{argument/DAG.md}.
\end{defnote}

%% ---- def-near-cluster --------------------------------------------------
\hypertarget{def:near-cluster}{}%
\begin{definition}[near-cluster\normalfont{} (\texttt{def-\allowbreak{}near-\allowbreak{}cluster})]\label{def:near-cluster}
\emph{Aliases:} rho-near deep cluster; near-deep cluster; cluster C(v); C(v); huddle cluster; cluster mass\par\smallskip
\textbf{Statement.} Fix an \hyperref[def:signed-idempotent]{exact signed idempotent} $P$ with nonempty
\hyperref[def:visible-set]{visible set} $\mathcal W$, a \hyperref[def:height]{hidden top vertex} $v$ of height $H$,
$\tau=\sqrt{\delta(P)}$, $\rho=4\tau$, and a width constant $a\ge4$. Write
$d_j=\operatorname{dist}_1(p_j,\operatorname{conv}\{p_w:w\in\mathcal W\})$.

The \textbf{near-cluster} (or $\rho$-near deep cluster) of $v$ is
\[C(v)\;:=\;\{\,j:\ \lVert p_j-p_v\rVert_1<4\tau\ \text{ and }\ d_j>a\,\tau\,\}\]
--- rows simultaneously $\rho$-near $v$ \textbf{and} deep (both conditions required). The \textbf{cluster mass}
is $\sum_{j\in C(v)}\max(P_{vj},0)$; a top is \textbf{heavy} when this is $\ge1-\theta_0$. A
\textbf{cluster vertex} is a geometrically distinct row vertex $u$ carrying positive disintegrated weight
from $C(v)$-rows (or, in the re-rooted assembly, $u:=v$ itself).

\textbf{Notes / provenance.} Project-original; the exact summand set of
\hyperlink{dag:conj-near-cluster-absorption}{\texttt{conj-\allowbreak{}near-\allowbreak{}cluster-\allowbreak{}absorption}} and the pinned target of the W54 huddle-charge decomposition tree.
The width $a$ is a calibration constant (any universal $a\ge4$; $a=16$ in the W54 assembly), not
load-bearing downstream. NOTE ON DIVERGENCE: the bare word ``cluster'' is also used generically for an
arbitrary row-index subset $C$ (e.g. \texttt{lem-\allowbreak{}cluster-\allowbreak{}return-\allowbreak{}flow}'s return flow over any subset) --- that
generic set-theoretic use is common knowledge and is NOT this term; this shard defines only the
huddle near-cluster $C(v)$. \texttt{status:\allowbreak{} draft} --- A+B sign-off pending (Rule 7). Related: \hyperref[def:co-top]{\texttt{def-\allowbreak{}co-\allowbreak{}top}},
\hyperref[def:slab]{\texttt{def-\allowbreak{}slab}}.
\end{definition}
\begin{defnote}
\textbf{Kind.} original (introduced by this project); \texttt{status: draft}.
\par
\textbf{Provenance.} \texttt{source: internal}; locus \texttt{internal;\allowbreak{} first pinned in argument/\allowbreak{}lemmas/\allowbreak{}conj-\allowbreak{}near-\allowbreak{}cluster-\allowbreak{}absorption.\allowbreak{}md}; no source hash (not a literature transcription); record: project-\allowbreak{}original;\allowbreak{} sign-\allowbreak{}off pending (Rule 7 lock).\allowbreak{} First pinned by conj-\allowbreak{}near-\allowbreak{}cluster-\allowbreak{}absorption (the summand set C(v)) and the W54 huddle-\allowbreak{}charge tree.
\par
\textbf{Used by} 3 registry results, none of them reproduced in this report (exploration track; see \texttt{argument/INDEX.md} and \texttt{report/UNWIRED.md}).
\par
\textbf{Canonical shard.} \texttt{definitions/\allowbreak{}def-\allowbreak{}near-\allowbreak{}cluster.\allowbreak{}md}; twins: \texttt{definitions/INDEX.md}, \texttt{argument/DAG.md}.
\end{defnote}

%% ---- def-near-positive-projection --------------------------------------
\hypertarget{def:near-positive-projection}{}%
\begin{definition}[near-positive projection\normalfont{} (\texttt{def-\allowbreak{}near-\allowbreak{}positive-\allowbreak{}projection})]\label{def:near-positive-projection}
\emph{Aliases:} $\delta$-positive unital idempotent; delta-positive unital idempotent; near-idempotent stochastic map; near-positive idempotent\par\smallskip
\textbf{Statement (commutative shadow).} A \emph{near-positive projection} is a linear map
$R\colon\ell^\infty_n\to\ell^\infty_n$ that is an exact unital idempotent, $R^2=R$ and
$R(\mathbf 1)=\mathbf 1$, and is \emph{$\delta$-positive}: for every $0\le x$ (entrywise) with
$\lVert x\rVert_\infty\le1$ one has $Rx\ge-\delta\,\mathbf 1$, for a small parameter
$\delta\in[0,\delta_0)$. As a matrix this is exactly an \hyperref[def:signed-idempotent]{exact signed idempotent}
whose \hyperref[def:negative-mass]{negative mass} satisfies $\delta(R)\le\delta$ (and $\lVert R\rVert\le1+O(\delta)$).

The motivating operator-algebra example is the spectral idempotent $P=\theta(2\Phi-\mathbf 1)$ built from
an \hyperref[def:almost-idempotent]{\texttt{def-\allowbreak{}almost-\allowbreak{}idempotent}} $\Phi$, with $\delta=O(\eta)$; such a $P$ need \textbf{not} be a genuinely
positive map. \emph{Near-positive-projection stability} (the open hypothesis parenting \texttt{op-\allowbreak{}classical}) asks
whether every near-positive projection is within $C\sqrt\delta$ (operator norm) of a genuine \emph{positive}
unital idempotent --- a \hyperref[def:stochastic]{stochastic idempotent} $E$; the $\sqrt{}$ exponent is sharp
(Hume's $3\times3$ family).

\textbf{Notes / provenance.} Adopted from \texttt{..\allowbreak{}/\allowbreak{}almost-\allowbreak{}idempotent-\allowbreak{}positive-\allowbreak{}maps/\allowbreak{}definitions/\allowbreak{}def-\allowbreak{}near-\allowbreak{}positive-\allowbreak{}projection.\allowbreak{}md}
(B's factorization program, \texttt{op:\allowbreak{}npps}), taken here as the commutative shadow / general parent of the
classical object. This is the map/positivity framing of the same object as \hyperref[def:signed-idempotent]{\texttt{def-\allowbreak{}signed-\allowbreak{}idempotent}}
(matrix/geometry framing): $R$ is $\delta$-positive iff $\delta(R)\le\delta$. The \emph{stability} statement is
an open registry problem, not part of this definition.
\end{definition}
\begin{defnote}
\textbf{Kind.} original (introduced by this project); \texttt{status: locked}.
\par
\textbf{Provenance.} \texttt{source: internal}; locus \texttt{adopted from ..\allowbreak{}/\allowbreak{}almost-\allowbreak{}idempotent-\allowbreak{}positive-\allowbreak{}maps/\allowbreak{}definitions/\allowbreak{}def-\allowbreak{}near-\allowbreak{}positive-\allowbreak{}projection.\allowbreak{}md (commutative shadow)}; no source hash (not a literature transcription); record: adopted from ..\allowbreak{}/\allowbreak{}almost-\allowbreak{}idempotent-\allowbreak{}positive-\allowbreak{}maps definitions/\allowbreak{}def-\allowbreak{}near-\allowbreak{}positive-\allowbreak{}projection.\allowbreak{}md (B's factorization program,\allowbreak{} op:\allowbreak{}npps);\allowbreak{} commutative shadow agreed for this repo.
\par
\textbf{Used by} 2 registry results, none of them reproduced in this report (exploration track; see \texttt{argument/INDEX.md} and \texttt{report/UNWIRED.md}).
\par
\textbf{Canonical shard.} \texttt{definitions/\allowbreak{}def-\allowbreak{}near-\allowbreak{}positive-\allowbreak{}projection.\allowbreak{}md}; twins: \texttt{definitions/INDEX.md}, \texttt{argument/DAG.md}.
\end{defnote}

%% ---- def-pivot ---------------------------------------------------------
\hypertarget{def:pivot}{}%
\begin{definition}[pivot (chart)\normalfont{} (\texttt{def-\allowbreak{}pivot})]\label{def:pivot}
\emph{Aliases:} pivot; maximal pivot; pivot-removing move; pivot-removing chart; theta-half chart; theta-half admissible; Phi-argmin chart; actual-row chart; Phi\_r; Psi\_j; Gamma\_j\par\smallskip
\textbf{Statement.} Fix a rank-$3$ \hyperref[def:signed-idempotent]{exact signed idempotent} $P$ with rows $p_i$.

\begin{itemize}
\item \textbf{Actual-row chart.} $U=(u_0,u_1,u_2)$, a triple of row indices whose rows $p_{u_0},p_{u_1},p_{u_2}$
form a basis of the row space. Old coordinates: $a_q(i)$ with $p_i=\sum_q a_q(i)\,p_{u_q}$.
\item \textbf{Charge functionals.} $E_r(i)=\max\!\big(\sum_{q\ne r}\max(-a_q(i),0)-(1-a_r(i)),\,0\big)$,
$\Phi_r(U)=\sum_i\max(P_{u_r i},0)\,E_r(i)$, and $\Phi(U)=\max_r\Phi_r(U)$.
\item \textbf{Volume gauge.} $\operatorname{Vol}(U)$ the Gram volume of the chart, $\operatorname{Vol}_{\max}(P)$
the maximum Gram volume over actual-row charts, $m_U=\operatorname{Vol}(U)/\operatorname{Vol}_{\max}(P)$.
$U$ is \textbf{theta-half} if $m_U\ge\tfrac12$; a \textbf{theta-half $\Phi$-argmin chart} is a theta-half $U$
minimizing $\Phi$ among theta-half actual-row charts.
\item \textbf{Pivot.} A \textbf{pivot} is a chart index $s$; a \textbf{maximal pivot} has $\Phi_s(U)=\Phi(U)$.
\item \textbf{Pivot-removing move.} For a non-chart row $j$ with $c=a_s(j)\ne0$, the \textbf{pivot-removing chart}
is $V_j=U-u_s+j$; it is \textbf{theta-half admissible} when $|a_s(j)|\,m_U\ge\tfrac12$ (volume factor
$\operatorname{Vol}(V_j)/\operatorname{Vol}(U)=|a_s(j)|$). New coordinates on $V_j$:
$a_s^{\,j}(i)=a_s(i)/c$ and $a_q^{\,j}(i)=a_q(i)-a_s(i)a_q(j)/c$ for $q\ne s$. The
\textbf{transferred charges} are
$\Psi_j=\sum_i\max(P_{ji},0)\,\max\!\big(\sum_{q\ne s}\max(-a_q^{\,j}(i),0)-(1-a_s^{\,j}(i)),0\big)$
and $\Gamma_j=\max_{r\ne s}\sum_i\max(P_{u_r i},0)\,\max\!\big(\sum_{q\ne r}\max(-a_q^{\,j}(i),0)-(1-a_r^{\,j}(i)),0\big)$.
\end{itemize}

Pivot-removing max-stationarity (\texttt{lem-\allowbreak{}pivot-\allowbreak{}removing-\allowbreak{}move}) is the bound $\Phi_s(U)\le\max(\Psi_j,\Gamma_j)$.

\textbf{Notes / provenance.} Project-original; the ``pivot'' here is the chart-swap sense (a basis row
exchanged for an external row), NOT the generic linear-algebra pivot. This shard is the single home of
the rank-$3$ chart/volume vocabulary that was inlined verbatim in the contracts of
\texttt{lem-\allowbreak{}pivot-\allowbreak{}removing-\allowbreak{}move}, \texttt{conj-\allowbreak{}b-\allowbreak{}restricted} and \texttt{conj-\allowbreak{}gamma-\allowbreak{}emptiness} (the latter two had
their notation preambles moved here 2026-07-10). \texttt{status:\allowbreak{} draft} --- A+B sign-off pending (Rule 7).
Related: \hyperref[def:signed-idempotent]{\texttt{def-\allowbreak{}signed-\allowbreak{}idempotent}}, \hyperref[def:negative-mass]{\texttt{def-\allowbreak{}negative-\allowbreak{}mass}}.
\end{definition}
\begin{defnote}
\textbf{Kind.} original (introduced by this project); \texttt{status: draft}.
\par
\textbf{Provenance.} \texttt{source: internal}; locus \texttt{internal;\allowbreak{} first pinned in argument/\allowbreak{}lemmas/\allowbreak{}lem-\allowbreak{}pivot-\allowbreak{}removing-\allowbreak{}move.\allowbreak{}md (and inlined verbatim in conj-\allowbreak{}b-\allowbreak{}restricted /\allowbreak{} conj-\allowbreak{}gamma-\allowbreak{}emptiness)}; no source hash (not a literature transcription); record: project-\allowbreak{}original;\allowbreak{} sign-\allowbreak{}off pending (Rule 7 lock).\allowbreak{} First pinned by lem-\allowbreak{}pivot-\allowbreak{}removing-\allowbreak{}move.
\par
\textbf{Used by} 2 registry results, none of them reproduced in this report (exploration track; see \texttt{argument/INDEX.md} and \texttt{report/UNWIRED.md}).
\par
\textbf{Canonical shard.} \texttt{definitions/\allowbreak{}def-\allowbreak{}pivot.\allowbreak{}md}; twins: \texttt{definitions/INDEX.md}, \texttt{argument/DAG.md}.
\end{defnote}

%% ---- def-top-support-functional ----------------------------------------
\hypertarget{def:top-support-functional}{}%
\begin{definition}[top support functional\normalfont{} (\texttt{def-\allowbreak{}top-\allowbreak{}support-\allowbreak{}functional})]\label{def:top-support-functional}
\emph{Aliases:} top support functional phi; Phi\_v; support functional at the top; averaged phi-bar\par\smallskip
\textbf{Statement.} Fix an \hyperref[def:signed-idempotent]{exact signed idempotent} $P$ with nonempty
\hyperref[def:visible-set]{visible set} $\mathcal W$, write $C_{\mathcal W}=\operatorname{conv}\{p_w:w\in\mathcal W\}$,
and let $v$ be a \hyperref[def:height]{hidden top vertex} of height $H$.

A \textbf{top support functional} at $v$ is an affine $\phi:\mathbb R^n\to\mathbb R$ with
\[\phi(p_v)=H,\qquad \phi\le 0\ \text{on}\ C_{\mathcal W},\qquad \phi\ \text{is $1$-Lipschitz for }\ell^1.\]
The set of all such is $\Phi_v$; it is nonempty (\hyperlink{dag:lem-top-deficit-price}{\texttt{lem-\allowbreak{}top-\allowbreak{}deficit-\allowbreak{}price}}), convex and compact, and
by \hyperlink{dag:lem-top-support-dual-face}{\texttt{lem-\allowbreak{}top-\allowbreak{}support-\allowbreak{}dual-\allowbreak{}face}} equals $\{\phi_y(x)=y\cdot x-h_{C}(y):y\in Y_v\}$ where
$h_C(y)=\sup_{c\in C_{\mathcal W}}y\cdot c$ and $Y_v=\{y:\lVert y\rVert_\infty\le1,\ y\cdot p_v-h_C(y)=H\}$.
A finite convex average of top support functionals (an \textbf{averaged $\bar\phi$}) is again a top support
functional (the three defining conditions are convex).

\textbf{Notes / provenance.} Project-original; the affine functional that realizes the height of the top
and pins the \hyperref[def:top-deficit]{top-deficit} $z_j=H-\phi(p_j)\ge0$. It is the charging channel of
\hyperlink{dag:lem-top-deficit-price}{\texttt{lem-\allowbreak{}top-\allowbreak{}deficit-\allowbreak{}price}}. \texttt{status:\allowbreak{} draft} --- A+B sign-off pending (Rule 7). Related:
\hyperref[def:top-deficit]{\texttt{def-\allowbreak{}top-\allowbreak{}deficit}}, \hyperref[def:height]{\texttt{def-\allowbreak{}height}}, \hyperref[def:visible-set]{\texttt{def-\allowbreak{}visible-\allowbreak{}set}}.
\end{definition}
\begin{defnote}
\textbf{Kind.} original (introduced by this project); \texttt{status: draft}.
\par
\textbf{Provenance.} \texttt{source: internal}; locus \texttt{internal;\allowbreak{} first pinned in argument/\allowbreak{}lemmas/\allowbreak{}lem-\allowbreak{}top-\allowbreak{}deficit-\allowbreak{}price.\allowbreak{}md and lem-\allowbreak{}top-\allowbreak{}support-\allowbreak{}dual-\allowbreak{}face.\allowbreak{}md}; no source hash (not a literature transcription); record: project-\allowbreak{}original;\allowbreak{} sign-\allowbreak{}off pending (Rule 7 lock).\allowbreak{} First pinned by lem-\allowbreak{}top-\allowbreak{}deficit-\allowbreak{}price (existence) + lem-\allowbreak{}top-\allowbreak{}support-\allowbreak{}dual-\allowbreak{}face (the dual face Phi\_v).
\par
\textbf{Used by} 39 registry results, none of them reproduced in this report (exploration track; see \texttt{argument/INDEX.md} and \texttt{report/UNWIRED.md}).
\par
\textbf{Canonical shard.} \texttt{definitions/\allowbreak{}def-\allowbreak{}top-\allowbreak{}support-\allowbreak{}functional.\allowbreak{}md}; twins: \texttt{definitions/INDEX.md}, \texttt{argument/DAG.md}.
\end{defnote}

%% ---- def-selected-corner -----------------------------------------------
\hypertarget{def:selected-corner}{}%
\begin{definition}[selected-corner configuration\normalfont{} (\texttt{def-\allowbreak{}selected-\allowbreak{}corner})]\label{def:selected-corner}
\emph{Aliases:} selected corner; corner configuration; disintegration kernel xi; Gamma\_f; M\_X; M\_I; M\_D; B\_F; B\_N; C\_f\par\smallskip
\textbf{Statement.} Fix a finite \hyperref[def:signed-idempotent]{exact signed idempotent} $P$ with $0<\delta(P)$,
$\tau=\sqrt{\delta(P)}$, $D=2+4\delta(P)$, nonempty \hyperref[def:visible-set]{visible set} $\mathcal W$,
$C_{\mathcal W}=\operatorname{conv}\{p_w:w\in\mathcal W\}$, and a \hyperref[def:height]{hidden top} $v$ of
height $H>16\tau$. A \textbf{selected-corner configuration} is such $(P,v)$ together with:

\begin{itemize}
\item an affine $\phi$ with $\phi(p_v)=H$, $\phi\le0$ on $C_{\mathcal W}$, and $|\phi(a)-\phi(b)|\le\lVert a-b\rVert_1$
(a \hyperref[def:top-support-functional]{top support functional}), an \hyperref[def:exposed]{admissible exposer} $h$ at
$v$, and a \textbf{selected corner row} $f$ with $\lVert p_f-p_v\rVert_1\ge4\tau$,
$\operatorname{dist}_1(p_f,C_{\mathcal W})>H-4\tau$ (\hyperref[def:co-top]{co-top} and $\rho$-far) satisfying
the corner inequality $2(H-\phi(p_f))/D+h(p_f)\le 12\tau/13$;
\item a \textbf{disintegration kernel} $\xi_x(u)$ from row points to geometrically distinct row vertices,
constant on clone fibers, with $p_x=\sum_u\xi_x(u)p_u$ and Dirac at vertex points.
\end{itemize}

Write $P^+_{fx}=\sum_{j:p_j=p_x}\max(P_{fj},0)$, $\Gamma_f(x,u)=P^+_{fx}\,\xi_x(u)$,
$C_f=\{(x,u):H-\phi(p_x)<4\tau,\ h(p_x)<4\tau,\ H-\phi(p_u)<4\tau,\ h(p_u)<4\tau\}$,
$B_F=C_f\cap\{\lVert p_u-p_v\rVert_1\ge4\tau\}$, $B_N=C_f\cap\{\lVert p_u-p_v\rVert_1<4\tau\}$. For a
block $B\in\{B_F,B_N\}$ the three \textbf{corner masses} are
\[M_X(B)=\Gamma_f\{(x,u)\in B:p_x\ne p_u\},\quad
  M_I(B)=\Gamma_f\{(x,u)\in B:p_x=p_u,\ K_T(u)\cap K_O(u)\ne\emptyset\},\]
\[M_D(B)=\Gamma_f\{(x,u)\in B:p_x=p_u,\ K_T(u)\cap K_O(u)=\emptyset\},\]
with $K_T(u),K_O(u)$ the \hyperref[def:actor-hull]{always-tight hulls} at $u$ (off-diagonal $X$ = distinct
$p_x\ne p_u$; intersection $I$ = diagonal with intersecting hulls; deep $D$ = diagonal with disjoint
hulls).

\textbf{Notes / provenance.} Project-original; the single home of the SL1a ``corner'' machinery inlined
verbatim in \hyperlink{dag:lem-sl1a-three-cell-reduction}{\texttt{lem-\allowbreak{}sl1a-\allowbreak{}three-\allowbreak{}cell-\allowbreak{}reduction}} and the three cell conjectures. The three cells partition
the selected corner's mass into the $X/I/D$ diagonal cells. \texttt{status:\allowbreak{} draft} --- A+B sign-off pending
(Rule 7). Related: \hyperref[def:actor-hull]{\texttt{def-\allowbreak{}actor-\allowbreak{}hull}}, \hyperref[def:co-top]{\texttt{def-\allowbreak{}co-\allowbreak{}top}}, \hyperref[def:top-support-functional]{\texttt{def-\allowbreak{}top-\allowbreak{}support-\allowbreak{}functional}}, \hyperref[def:top-deficit]{\texttt{def-\allowbreak{}top-\allowbreak{}deficit}}.
\end{definition}
\begin{defnote}
\textbf{Kind.} original (introduced by this project); \texttt{status: draft}.
\par
\textbf{Provenance.} \texttt{source: internal}; locus \texttt{internal;\allowbreak{} first pinned (inlined verbatim) in argument/\allowbreak{}lemmas/\allowbreak{}lem-\allowbreak{}sl1a-\allowbreak{}three-\allowbreak{}cell-\allowbreak{}reduction.\allowbreak{}md and the three conj-\allowbreak{}sl1a-\allowbreak{}*-\allowbreak{}cell shards}; no source hash (not a literature transcription); record: project-\allowbreak{}original;\allowbreak{} sign-\allowbreak{}off pending (Rule 7 lock).\allowbreak{} First pinned by lem-\allowbreak{}sl1a-\allowbreak{}three-\allowbreak{}cell-\allowbreak{}reduction (the SL1a three-\allowbreak{}cell surface).
\par
\textbf{Used by} 37 registry results, none of them reproduced in this report (exploration track; see \texttt{argument/INDEX.md} and \texttt{report/UNWIRED.md}).
\par
\textbf{Canonical shard.} \texttt{definitions/\allowbreak{}def-\allowbreak{}selected-\allowbreak{}corner.\allowbreak{}md}; twins: \texttt{definitions/INDEX.md}, \texttt{argument/DAG.md}.
\end{defnote}

%% ---- def-top-deficit ---------------------------------------------------
\hypertarget{def:top-deficit}{}%
\begin{definition}[top-deficit\normalfont{} (\texttt{def-\allowbreak{}top-\allowbreak{}deficit})]\label{def:top-deficit}
\emph{Aliases:} top deficit; z\_j; deficit; top-deficit slab; Z\_v(f); top-deficit supremum\par\smallskip
\textbf{Statement.} In the setting of \hyperref[def:top-support-functional]{\texttt{def-\allowbreak{}top-\allowbreak{}support-\allowbreak{}functional}} --- an
\hyperref[def:signed-idempotent]{exact signed idempotent} $P$, nonempty \hyperref[def:visible-set]{visible set},
\hyperref[def:height]{hidden top} $v$ of height $H$, and a top support functional $\phi\in\Phi_v$ --- the
\textbf{top-deficit} of a row $j$ under $\phi$ is
\[z_j\;:=\;H-\phi(p_j)\;\ge\;0,\]
nonnegative because $\phi(p_j)\le\operatorname{dist}_1(p_j,C_{\mathcal W})\le H$ (and bounded above,
$z_j\le 2+4\delta$). The \textbf{top-deficit supremum} of a row $f$ is
$Z_v(f):=\sup_{\phi\in\Phi_v}\big(H-\phi(p_f)\big)=\sup_{y\in Y_v}y\cdot(p_v-p_f)$
(\hyperlink{dag:lem-top-support-dual-face}{\texttt{lem-\allowbreak{}top-\allowbreak{}support-\allowbreak{}dual-\allowbreak{}face}}). A \textbf{top-deficit slab} is a band of rows cut by a threshold on $z$,
e.g. $\{j:z_j\ge L\}$.

\textbf{Notes / provenance.} Project-original; ``deficit'' alone always means this top-deficit in the
huddle-charge shards. The charging identity (\hyperlink{dag:lem-top-deficit-price}{\texttt{lem-\allowbreak{}top-\allowbreak{}deficit-\allowbreak{}price}}, from $P^2=P$ and $P\mathbf1=\mathbf1$)
is $\sum_j a_j^+z_j=\sum_j a_j^-z_j\le\nu_v(2+4\delta)$ with $a_j=P_{vj}$; its structural blind spot is
that rows in the $\rho$-ball of $v$ carry $z_j<4\tau$. \texttt{status:\allowbreak{} draft} --- A+B sign-off pending (Rule 7).
Related: \hyperref[def:top-support-functional]{\texttt{def-\allowbreak{}top-\allowbreak{}support-\allowbreak{}functional}}, \hyperref[def:slab]{\texttt{def-\allowbreak{}slab}}, \hyperref[def:co-top]{\texttt{def-\allowbreak{}co-\allowbreak{}top}}.
\end{definition}
\begin{defnote}
\textbf{Kind.} original (introduced by this project); \texttt{status: draft}.
\par
\textbf{Provenance.} \texttt{source: internal}; locus \texttt{internal;\allowbreak{} first pinned in argument/\allowbreak{}lemmas/\allowbreak{}lem-\allowbreak{}top-\allowbreak{}deficit-\allowbreak{}price.\allowbreak{}md and lem-\allowbreak{}top-\allowbreak{}support-\allowbreak{}dual-\allowbreak{}face.\allowbreak{}md}; no source hash (not a literature transcription); record: project-\allowbreak{}original;\allowbreak{} sign-\allowbreak{}off pending (Rule 7 lock).\allowbreak{} First pinned by lem-\allowbreak{}top-\allowbreak{}deficit-\allowbreak{}price (z\_j) + lem-\allowbreak{}top-\allowbreak{}support-\allowbreak{}dual-\allowbreak{}face (the supremum Z\_v(f)).
\par
\textbf{Used by} no registry result yet (vocabulary registered ahead of the argument).
\par
\textbf{Canonical shard.} \texttt{definitions/\allowbreak{}def-\allowbreak{}top-\allowbreak{}deficit.\allowbreak{}md}; twins: \texttt{definitions/INDEX.md}, \texttt{argument/DAG.md}.
\end{defnote}

%% ---- def-slab ----------------------------------------------------------
\hypertarget{def:slab}{}%
\begin{definition}[slab\normalfont{} (\texttt{def-\allowbreak{}slab})]\label{def:slab}
\emph{Aliases:} low slab; low-slab; deep low slab; very-low slab; top slab; top-slab; far slab; slab leakage\par\smallskip
\textbf{Statement.} Fix an \hyperref[def:signed-idempotent]{exact signed idempotent} $P$, $\tau=\sqrt{\delta(P)}$,
nonempty \hyperref[def:visible-set]{visible set} $\mathcal W$. A \textbf{slab} is the set of rows on which a fixed
affine functional lies in a prescribed threshold band. Two families recur, kept distinct here:

\begin{itemize}
\item \textbf{Low slab} (exposer-value): for an \hyperref[def:exposed]{optimal exposer} $h^*$ at a vertex and a
threshold $s$, the set $\{j:h^*(p_j)<s\}$. The \textbf{deep / very-low slab} takes $s=\kappa=\tau/4$ or
$s=\tau/8$; combined with depth ($d_j>a\tau$) it is the summand set of \hyperlink{dag:conj-low-slab-cap}{\texttt{conj-\allowbreak{}low-\allowbreak{}slab-\allowbreak{}cap}} and
\hyperlink{dag:conj-far-low-slab-cap}{\texttt{conj-\allowbreak{}far-\allowbreak{}low-\allowbreak{}slab-\allowbreak{}cap}}. The pincer \hyperlink{dag:lem-cs-low-slab-pincer}{\texttt{lem-\allowbreak{}cs-\allowbreak{}low-\allowbreak{}slab-\allowbreak{}pincer}} controls its complementary shell
$\{s\le h^*<\kappa\}$.
\item \textbf{Top slab} (depth / \hyperref[def:top-deficit]{top-deficit}): the near-top band
$\{j:d_j>H-c\tau\}$ (equivalently small top-deficit $z_j<c'\tau$); the \textbf{far slab} additionally
imposes $\rho$-farness. \texttt{lem-\allowbreak{}top-\allowbreak{}slab-\allowbreak{}companion} forces a $\rho$-far row in the top slab.
\end{itemize}

\textbf{Notes / provenance.} Project-original; ``slab'' is overloaded in the shards between the
exposer-value cut ($h^*$) and the depth cut ($d$ / $z$) --- this shard records BOTH canonical instances
and which functional each uses, so downstream statements can name the functional explicitly and avoid
drift. \texttt{status:\allowbreak{} draft} --- A+B sign-off pending (Rule 7). Related: \hyperref[def:top-deficit]{\texttt{def-\allowbreak{}top-\allowbreak{}deficit}},
\hyperref[def:near-cluster]{\texttt{def-\allowbreak{}near-\allowbreak{}cluster}}, \hyperref[def:exposed]{\texttt{def-\allowbreak{}exposed}}.
\end{definition}
\begin{defnote}
\textbf{Kind.} original (introduced by this project); \texttt{status: draft}.
\par
\textbf{Provenance.} \texttt{source: internal}; locus \texttt{internal;\allowbreak{} first pinned across argument/\allowbreak{}lemmas/\allowbreak{}conj-\allowbreak{}low-\allowbreak{}slab-\allowbreak{}cap.\allowbreak{}md,\allowbreak{} lem-\allowbreak{}cs-\allowbreak{}low-\allowbreak{}slab-\allowbreak{}pincer.\allowbreak{}md,\allowbreak{} lem-\allowbreak{}top-\allowbreak{}slab-\allowbreak{}companion.\allowbreak{}md}; no source hash (not a literature transcription); record: project-\allowbreak{}original;\allowbreak{} sign-\allowbreak{}off pending (Rule 7 lock).\allowbreak{} First pinned by conj-\allowbreak{}low-\allowbreak{}slab-\allowbreak{}cap (exposer-\allowbreak{}value slab) and lem-\allowbreak{}top-\allowbreak{}slab-\allowbreak{}companion (depth slab).
\par
\textbf{Used by} 2 registry results, none of them reproduced in this report (exploration track; see \texttt{argument/INDEX.md} and \texttt{report/UNWIRED.md}).
\par
\textbf{Canonical shard.} \texttt{definitions/\allowbreak{}def-\allowbreak{}slab.\allowbreak{}md}; twins: \texttt{definitions/INDEX.md}, \texttt{argument/DAG.md}.
\end{defnote}


% The approximate $C^{*}$-algebra vocabulary (transcribed from the pinned source)
% GENERATED by scripts/gen-report-defs.py — DO NOT HAND-EDIT.
% Single source of truth: definitions/<id>.md (CLAUDE.md L2).  Re-render with
%   python3 scripts/gen-report-defs.py
% The gate `python3 scripts/gen-report-defs.py --check` fails when this file is stale.

% Layer file: The approximate $C^{*}$-algebra vocabulary (transcribed from the pinned source)
\subsection*{The approximate $C^{*}$-algebra vocabulary (transcribed from the pinned source)}

%% ---- def-column-hilbert-corner -----------------------------------------
\hypertarget{def:column-hilbert-corner}{}%
\begin{definition}[level-one and amplified column-Hilbert corner\normalfont{} (\texttt{def-\allowbreak{}column-\allowbreak{}hilbert-\allowbreak{}corner})]\label{def:column-hilbert-corner}
\emph{Aliases:} column-Hilbert corner; amplified column space\par\smallskip
\textbf{Byte-verbatim source text.}

\begin{defsourcequote}
\defsourceline{\begin{Lemma}\label{lem_PQ_Hilb}}
\defsourceline{Let $P$ and $Q$ be $\delta$-projections in an $\eps$-$C^*$ algebra. If $Q$ is one-dimensional, then $\calS_{P,Q}$ is a Hilbert space with the inner product $\braket{\cdot}{\cdot}$ defined by the equation}
\defsourceline{\begin{equation}\label{1dQ_ip}}
\defsourceline{Y^\dag\cdot X=\braket{Y}{X}\,\wt{Q}\qquad (X,Y\in\calS_{P,Q}),}
\defsourceline{\end{equation}}
\defsourceline{where $Y^\dag\cdot X=\Co_{Q}(Y^{\dag}X)$ and $\wt{Q}=\Co_Q(Q)$.}
\end{defsourcequote}

\begin{defsourcequote}
\defsourceline{\noindent\textbf{One-dimensional projections and the corresponding Hilbert spaces.} Lemma~\ref{lem_PQ_Hilb} says that if $P,Q\in\calA$ are $\delta$-projections and $Q$ is one-dimensional, then the space $\calS_{P,Q}$ is equipped with the Hermitian inner product given by \eqref{1dQ_ip}. The extended version of this space, $\CC^n\otimes\calS_{P,Q}=\Ma{n,1}\otimes\calS_{P,Q}$, satisfies essentially the same equation:}
\defsourceline{\begin{equation}\label{1dQ_ip_ext}}
\defsourceline{Y^\dag\cdot X=\braket{Y}{X}\,\wt{Q}\qquad (X,Y\in\Ma{n,1}\otimes\calS_{P,Q}),}
\defsourceline{\end{equation}}
\defsourceline{where $\braket{Y}{X}=\sum_{l}\braket{[Y]_{l1}}{[X]_{l1}}$.}
\end{defsourcequote}

\textbf{Status.} Draft transcription; ratification is required before locking. The
known false unsquared display at source lines 1551-1555 is deliberately not
part of this definition.

\textbf{Provenance.} Byte-verbatim from the pinned source at the recorded loci
and SHA256 prefix; the exclusion is required by \texttt{DESIGN-\allowbreak{}FUDW-\allowbreak{}DECOMP-\allowbreak{}v3.\allowbreak{}md}
\S{}4.1 and is consistent with \texttt{VERDICT-\allowbreak{}W74F-\allowbreak{}H-\allowbreak{}STAGE1.\allowbreak{}md}.
\end{definition}
\begin{defnote}
\textbf{Kind.} cited (byte-matched to a pinned source under \texttt{refs/}); \texttt{status: locked}.
\par
\textbf{Provenance.} \texttt{source: kitaev-\allowbreak{}2405.\allowbreak{}02434}; locus \texttt{approximate\_algebras.\allowbreak{}tex:\allowbreak{}1123-\allowbreak{}1128,\allowbreak{}1546-\allowbreak{}1550}; \textsc{sha256} prefix \texttt{e7eb512a2ec2438d}; record: user-\allowbreak{}ratified 2026-\allowbreak{}07-\allowbreak{}24 (tobiasosborne,\allowbreak{} in-\allowbreak{}session sign-\allowbreak{}off:\allowbreak{} byte-\allowbreak{}match-\allowbreak{}to-\allowbreak{}source criterion for cited;\allowbreak{} delegated ratification for consensus/\allowbreak{}original).
\par
\textbf{Used in this report by} \hyperref[lem:compcb-row-column-product]{\texttt{lem-\allowbreak{}compcb-\allowbreak{}row-\allowbreak{}column-\allowbreak{}product}}, \hyperref[lem:hcb-column-hilbert-squared]{\texttt{lem-\allowbreak{}hcb-\allowbreak{}column-\allowbreak{}hilbert-\allowbreak{}squared}} (2 registry results import it in total; see \texttt{argument/INDEX.md}).
\par
\textbf{Canonical shard.} \texttt{definitions/\allowbreak{}def-\allowbreak{}column-\allowbreak{}hilbert-\allowbreak{}corner.\allowbreak{}md}; twins: \texttt{definitions/INDEX.md}, \texttt{argument/DAG.md}.
\end{defnote}

%% ---- def-compressed-corner ---------------------------------------------
\hypertarget{def:compressed-corner}{}%
\begin{definition}[compressed corner\normalfont{} (\texttt{def-\allowbreak{}compressed-\allowbreak{}corner})]\label{def:compressed-corner}
\emph{Aliases:} compression map; compressed product; compressed unit\par\smallskip
\textbf{Byte-verbatim source text.}

\begin{defsourcequote}
\defsourceline{To each pair of projections $(P,Q)$ in a $C^*$ algebra $\calB$, one assigns the vector space $\calS_{P,Q}=\{PXQ:\,X\in\calB\}$. In this section, we discuss an approximate version of this construction. Recall that a $\delta$-projection in an $\eps$-$C^*$ algebra $\calA$ is a Hermitian element $P$ such that $\|P^2-P\|\le\delta$. For each pair $(P,Q)$ of $\delta$-projections, there are two slightly different maps that could be called ``compressions'': $\La_P\Ra_Q\colon X\mapsto P(XQ)$ and $\Ra_Q\La_P\colon X\mapsto (PX)Q$. We will use their symmetric combination, $\frac{1}{2}(\La_P\Ra_Q+\Ra_Q\La_P)$. It is an $O(\delta+\eps)$-idempotent element of the algebra of operators acting on $\calA$, and therefore, can be approximated by an idempotent using Proposition~\ref{prop_P}. The \emph{compression map} thus defined,}
\defsourceline{\begin{equation}}
\defsourceline{\Co_{P,Q}\colon\calA\to\calA,\qquad}
\defsourceline{\Co_{P,Q}=\theta(\La_P\Ra_Q+\Ra_Q\La_P-1),}
\defsourceline{\end{equation}}
\defsourceline{satisfies the equations}
\defsourceline{\begin{equation}}
\defsourceline{\Co_{P,Q}^2=\Co_{P,Q},\qquad\:}
\defsourceline{\Co_{P,Q}(X)^\dag=\Co_{Q,P}(X^\dag)\quad\: (X\in\calA),}
\defsourceline{\end{equation}}
\defsourceline{and both $\|\La_P\Ra_Q-\Co_{P,Q}\|$ and $\|\Ra_Q\La_P-\Co_{P,Q}\|$ are bounded by $O(\delta+\eps)$. The image of this map, $\calS_{P,Q}=\Img\Co_{P,Q}=\Ker(1-\Co_{P,Q})$, is a closed linear subspace of $\calA$. There is a variant of this construction where only $\La_P$ or only $\Ra_Q$ is applied, and one writes $\Co_{P,1},\calS_{P,1}$ or $\Co_{1,Q},\calS_{1,Q}$, respectively. (This is the same as $\Co_{P,I}$, etc.\ if the unit is exact.)}
\defsourcegap
\defsourceline{When $P=Q$, the abbreviations $\Co_{P}$, $\calS_P$ are used.}
\end{defsourcequote}

\begin{defsourcequote}
\defsourceline{For any triple $(P,Q,R)$ of $\delta$-projections, one defines the \emph{compressed product} between elements of the corresponding subspaces:}
\defsourceline{\begin{equation}\label{compr_prod}}
\defsourceline{(X,Y)\mapsto X\cdot Y\,\colon\,}
\defsourceline{\calS_{P,Q}\times\calS_{Q,R}\to \calS_{P,R},\qquad X\cdot Y=\Co_{P,R}(XY).}
\defsourceline{\end{equation}}
\defsourceline{It is close to the ambient product in $\calA$, namely, $\|X\cdot Y-XY\|\le O(\delta+\eps)\|X\|\ts\|Y\|$. If $P$ is a nonvanishing projection, the compressed product turns the Banach space $\calS_P$ (which is closed under the involution $X\mapsto X^\dag$) into an $O(\delta+\eps)$-$C^*$ algebra with unit $\wt{P}=\Co_{P}(P)$.\medskip}
\end{defsourcequote}

\textbf{Status.} Draft transcription; ratification is required before locking.

\textbf{Provenance.} Byte-verbatim from the pinned source at the recorded loci
and SHA256 prefix; provisioned by \texttt{DESIGN-\allowbreak{}FUDW-\allowbreak{}DECOMP-\allowbreak{}v3.\allowbreak{}md} \S{}4.1 and admitted
only for the safe subset by \texttt{VERDICT-\allowbreak{}FUDW-\allowbreak{}DECOMP-\allowbreak{}V3.\allowbreak{}md} \S{}D.
\end{definition}
\begin{defnote}
\textbf{Kind.} cited (byte-matched to a pinned source under \texttt{refs/}); \texttt{status: locked}.
\par
\textbf{Provenance.} \texttt{source: kitaev-\allowbreak{}2405.\allowbreak{}02434}; locus \texttt{approximate\_algebras.\allowbreak{}tex:\allowbreak{}1054-\allowbreak{}1066,\allowbreak{}1077-\allowbreak{}1082}; \textsc{sha256} prefix \texttt{e7eb512a2ec2438d}; record: user-\allowbreak{}ratified 2026-\allowbreak{}07-\allowbreak{}24 (tobiasosborne,\allowbreak{} in-\allowbreak{}session sign-\allowbreak{}off:\allowbreak{} byte-\allowbreak{}match-\allowbreak{}to-\allowbreak{}source criterion for cited;\allowbreak{} delegated ratification for consensus/\allowbreak{}original).
\par
\textbf{Used in this report by} \hyperref[lem:compcb-amplified-almost-containment]{\texttt{lem-\allowbreak{}compcb-\allowbreak{}amplified-\allowbreak{}almost-\allowbreak{}containment}}, \hyperref[lem:compcb-amplified-compression]{\texttt{lem-\allowbreak{}compcb-\allowbreak{}amplified-\allowbreak{}compression}}, \hyperref[lem:compcb-amplified-compression-identities]{\texttt{lem-\allowbreak{}compcb-\allowbreak{}amplified-\allowbreak{}compression-\allowbreak{}identities}}, \hyperref[lem:compcb-compressed-unit-action]{\texttt{lem-\allowbreak{}compcb-\allowbreak{}compressed-\allowbreak{}unit-\allowbreak{}action}}, \hyperref[lem:compcb-compressed-unit-norm]{\texttt{lem-\allowbreak{}compcb-\allowbreak{}compressed-\allowbreak{}unit-\allowbreak{}norm}}, \hyperref[lem:compcb-corner-algebra]{\texttt{lem-\allowbreak{}compcb-\allowbreak{}corner-\allowbreak{}algebra}} and 11 further reproduced results (17 registry results import it in total; see \texttt{argument/INDEX.md}).
\par
\textbf{Canonical shard.} \texttt{definitions/\allowbreak{}def-\allowbreak{}compressed-\allowbreak{}corner.\allowbreak{}md}; twins: \texttt{definitions/INDEX.md}, \texttt{argument/DAG.md}.
\end{defnote}

%% ---- def-delta-projection ----------------------------------------------
\hypertarget{def:delta-projection}{}%
\begin{definition}[$\delta$-projection and nonvanishing alternative\normalfont{} (\texttt{def-\allowbreak{}delta-\allowbreak{}projection})]\label{def:delta-projection}
\emph{Aliases:} delta-projection; nonvanishing delta-projection\par\smallskip
\textbf{Byte-verbatim source text.}

\begin{defsourcequote}
\defsourceline{A \emph{projection} in an $\eps$-$C^*$ algebra $\calA$ is a Hermitian element $P$ such that $P^2=P$. Such elements are generally hard to come by, see Remark~\ref{rem_X2}, so will content ourselves with \emph{$\delta$-projections} for sufficiently small $\delta$. They are defined by the conditions}
\defsourceline{\begin{equation}\label{delta_P}}
\defsourceline{P^\dag=P,\qquad \|P^2-P\|\le\delta.}
\defsourceline{\end{equation}}
\end{defsourcequote}

\begin{defsourcequote}
\defsourceline{\begin{equation}\label{P_alternatives}}
\defsourceline{\|P\|\le O(\delta)\quad \text{or}\quad \bigl|\|P\|-1\bigr|\le O(\delta+\eps).}
\defsourceline{\end{equation}}
\defsourceline{A $\delta$-projection is called \emph{nonvanishing} if the second alternative holds. If $P$ is a $\delta$-projection, then $I-P$ is a $\delta'$-projection, where $\delta'=\delta$ if $\calA$ has exact unit and $\delta'=\delta+O(\eps)$ in general. A $\delta$-projection $P$ is called \emph{nontrivial} if both $P$ and $I-P$ are nonvanishing.}
\end{defsourcequote}

\textbf{Status.} Draft transcription; ratification is required before locking.

\textbf{Provenance.} Byte-verbatim from the pinned source at the recorded loci
and SHA256 prefix; provisioned by \texttt{DESIGN-\allowbreak{}FUDW-\allowbreak{}DECOMP-\allowbreak{}v3.\allowbreak{}md} \S{}4.1 and admitted
only for the safe subset by \texttt{VERDICT-\allowbreak{}FUDW-\allowbreak{}DECOMP-\allowbreak{}V3.\allowbreak{}md} \S{}D.
\end{definition}
\begin{defnote}
\textbf{Kind.} cited (byte-matched to a pinned source under \texttt{refs/}); \texttt{status: locked}.
\par
\textbf{Provenance.} \texttt{source: kitaev-\allowbreak{}2405.\allowbreak{}02434}; locus \texttt{approximate\_algebras.\allowbreak{}tex:\allowbreak{}917-\allowbreak{}920,\allowbreak{}926-\allowbreak{}929}; \textsc{sha256} prefix \texttt{e7eb512a2ec2438d}; record: user-\allowbreak{}ratified 2026-\allowbreak{}07-\allowbreak{}24 (tobiasosborne,\allowbreak{} in-\allowbreak{}session sign-\allowbreak{}off:\allowbreak{} byte-\allowbreak{}match-\allowbreak{}to-\allowbreak{}source criterion for cited;\allowbreak{} delegated ratification for consensus/\allowbreak{}original).
\par
\textbf{Used in this report by} \hyperref[conj:extcb]{\texttt{conj-\allowbreak{}extcb}}, \hyperref[conj:hcb]{\texttt{conj-\allowbreak{}hcb}}, \hyperref[lem:compcb-amplified-almost-containment]{\texttt{lem-\allowbreak{}compcb-\allowbreak{}amplified-\allowbreak{}almost-\allowbreak{}containment}}, \hyperref[lem:compcb-amplified-compression]{\texttt{lem-\allowbreak{}compcb-\allowbreak{}amplified-\allowbreak{}compression}}, \hyperref[lem:compcb-amplified-compression-identities]{\texttt{lem-\allowbreak{}compcb-\allowbreak{}amplified-\allowbreak{}compression-\allowbreak{}identities}}, \hyperref[lem:compcb-compressed-unit-action]{\texttt{lem-\allowbreak{}compcb-\allowbreak{}compressed-\allowbreak{}unit-\allowbreak{}action}} and 9 further reproduced results (15 registry results import it in total; see \texttt{argument/INDEX.md}).
\par
\textbf{Canonical shard.} \texttt{definitions/\allowbreak{}def-\allowbreak{}delta-\allowbreak{}projection.\allowbreak{}md}; twins: \texttt{definitions/INDEX.md}, \texttt{argument/DAG.md}.
\end{defnote}

%% ---- def-epsilon-cstar-algebra -----------------------------------------
\hypertarget{def:epsilon-cstar-algebra}{}%
\begin{definition}[$\varepsilon$-$C^{*}$-algebra\normalfont{} (\texttt{def-\allowbreak{}epsilon-\allowbreak{}cstar-\allowbreak{}algebra})]\label{def:epsilon-cstar-algebra}
\emph{Aliases:} epsilon-$C^{*}$-algebra; approximate $C^{*}$-algebra\par\smallskip
\textbf{Byte-verbatim source text.}

\begin{defsourcequote}
\defsourceline{\begin{Definition}}
\defsourceline{An \emph{$\eps$-Banach algebra} is a Banach space $\calA$ endowed with a bilinear multiplication map $\calA\times \calA\to \calA$ such that}
\defsourceline{\begin{alignat}{2}}
\defsourceline{\label{ax_prodnorm}}
\defsourceline{\|XY\| &\le (1+\eps)\ts\|X\|\ts\|Y\|\qquad &&(X,Y\in \calA),\\[2pt]}
\defsourceline{\label{ax_assoc}}
\defsourceline{\|(XY)Z-X(YZ)\| &\le \eps\ts\|X\|\ts\|Y\|\ts\|Z\|\qquad &&(X,Y,Z\in \calA).}
\defsourceline{\end{alignat}}
\defsourceline{Such an algebra is called \emph{$\eps'$-commutative} if}
\defsourceline{\begin{equation}}
\defsourceline{\label{ax_comm}}
\defsourceline{\|XY-YX\|\le \eps'\ts\|X\|\ts\|Y\|\qquad (X,Y\in \calA).}
\defsourceline{\end{equation}}
\defsourceline{A \emph{$*\eps$-Banach algebra} is a complex $\eps$-Banach algebra with a conjugate linear involution $X\mapsto X^\dag$ satisfying the equations}
\defsourceline{\begin{equation}}
\defsourceline{\label{ax_*}}
\defsourceline{\|X^\dag\|=\|X\|,\qquad (XY)^\dag=Y^\dag X^\dag\qquad (X,Y\in \calA).}
\defsourceline{\end{equation}}
\defsourceline{An \emph{$\eps$-$C^*$ algebra} is one with this additional property:}
\defsourceline{\begin{equation}}
\defsourceline{\label{ax_C*}}
\defsourceline{\|X^{\dag}X\|\ge (1-\eps)\ts\|X\|^{2}\qquad (X\in \calA).}
\defsourceline{\end{equation}}
\defsourceline{(A bound from the other side, $\|X^{\dag}X\|\le (1+\eps)\ts\|X\|^{2}$, follows from \eqref{ax_prodnorm} and \eqref{ax_*}.) We assume that all algebras are unital. The unit element $I\in\calA$ should satisfy the approximate or exact conditions}
\defsourceline{\begin{alignat}{4}}
\defsourceline{\label{ax_eps_unit}}
\defsourceline{\|XI&-X\|\le\eps\ts\|X\|,\qquad &\|IX&-X\|\le\eps\ts\|X\|,\qquad}
\defsourceline{&\bigl|\|I\|&-1\bigr|\le\eps\qquad &&\text{(by default)},\quad \text{or}\\[2pt]}
\defsourceline{\label{ax_exact_unit}}
\defsourceline{XI&=X,\qquad &IX&=X,\qquad}
\defsourceline{&\|I\|&=1\qquad &&\text{(if specified)}.}
\defsourceline{\end{alignat}}
\defsourceline{If the involution is defined, one also requires that $I^\dag=I$.}
\defsourceline{\end{Definition}}
\end{defsourcequote}

\textbf{Status.} Draft transcription; ratification is required before locking.

\textbf{Provenance.} Byte-verbatim from the pinned source at the recorded locus
and SHA256 prefix; provisioned by \texttt{DESIGN-\allowbreak{}FUDW-\allowbreak{}DECOMP-\allowbreak{}v3.\allowbreak{}md} \S{}4.1 and admitted
only for the safe subset by \texttt{VERDICT-\allowbreak{}FUDW-\allowbreak{}DECOMP-\allowbreak{}V3.\allowbreak{}md} \S{}D.
\end{definition}
\begin{defnote}
\textbf{Kind.} cited (byte-matched to a pinned source under \texttt{refs/}); \texttt{status: locked}.
\par
\textbf{Provenance.} \texttt{source: kitaev-\allowbreak{}2405.\allowbreak{}02434}; locus \texttt{approximate\_algebras.\allowbreak{}tex:\allowbreak{}407-\allowbreak{}440}; \textsc{sha256} prefix \texttt{e7eb512a2ec2438d}; record: user-\allowbreak{}ratified 2026-\allowbreak{}07-\allowbreak{}24 (tobiasosborne,\allowbreak{} in-\allowbreak{}session sign-\allowbreak{}off:\allowbreak{} byte-\allowbreak{}match-\allowbreak{}to-\allowbreak{}source criterion for cited;\allowbreak{} delegated ratification for consensus/\allowbreak{}original).
\par
\textbf{Used in this report by} \hyperref[lem:extcb-corner-dimension-additivity]{\texttt{lem-\allowbreak{}extcb-\allowbreak{}corner-\allowbreak{}dimension-\allowbreak{}additivity}}, \hyperref[lem:extcb-one-dimensional-corner-dimension]{\texttt{lem-\allowbreak{}extcb-\allowbreak{}one-\allowbreak{}dimensional-\allowbreak{}corner-\allowbreak{}dimension}}, \hyperref[lem:extcb-one-dimensional-product]{\texttt{lem-\allowbreak{}extcb-\allowbreak{}one-\allowbreak{}dimensional-\allowbreak{}product}}, \hyperref[lem:stage1-exact-unit-rectification]{\texttt{lem-\allowbreak{}stage1-\allowbreak{}exact-\allowbreak{}unit-\allowbreak{}rectification}} (4 registry results import it in total; see \texttt{argument/INDEX.md}).
\par
\textbf{Canonical shard.} \texttt{definitions/\allowbreak{}def-\allowbreak{}epsilon-\allowbreak{}cstar-\allowbreak{}algebra.\allowbreak{}md}; twins: \texttt{definitions/INDEX.md}, \texttt{argument/DAG.md}.
\end{defnote}

%% ---- def-extended-epsilon-cstar-algebra --------------------------------
\hypertarget{def:extended-epsilon-cstar-algebra}{}%
\begin{definition}[extended $\varepsilon$-$C^{*}$-algebra\normalfont{} (\texttt{def-\allowbreak{}extended-\allowbreak{}epsilon-\allowbreak{}cstar-\allowbreak{}algebra})]\label{def:extended-epsilon-cstar-algebra}
\emph{Aliases:} extended epsilon-$C^{*}$-algebra; extended approximate $C^{*}$-algebra\par\smallskip
\textbf{Byte-verbatim source text.}

\begin{defsourcequote}
\defsourceline{An \emph{extended $\eps$-$C^*$ algebra} is a complete self-adjoint operator space $\calA$ with a multiplication and a unit that make each space $\Ma{n}\otimes\calA$ into an $\eps$-$C^*$ algebra. An \emph{extended $\delta$-homomorphism} is a linear map $v\colon\calA'\to\calA''$ such that for each $n$, the map $1_{\Ma{n}}\otimes v\colon \Ma{n}\otimes\calA'\to\Ma{n}\otimes\calA''$ is a $\delta$-homomorphism.}
\end{defsourcequote}

\textbf{Notation.} Registry contracts write the source's \texttt{\textbackslash{}(\textbackslash{}eps\textbackslash{})} as
\texttt{\textbackslash{}(\textbackslash{}varepsilon\textbackslash{})} and \texttt{\textbackslash{}(\textbackslash{}Ma n\textbackslash{})} as \texttt{\textbackslash{}(M\_n\textbackslash{})}.

\textbf{Provenance.} Byte-verbatim from the pinned source at the recorded locus and
SHA256 prefix.  User-ratified and locked 2026-07-24 (tobiasosborne, in-session sign-off).
\end{definition}
\begin{defnote}
\textbf{Kind.} cited (byte-matched to a pinned source under \texttt{refs/}); \texttt{status: locked}.
\par
\textbf{Provenance.} \texttt{source: kitaev-\allowbreak{}2405.\allowbreak{}02434}; locus \texttt{approximate\_algebras.\allowbreak{}tex:\allowbreak{}1477-\allowbreak{}1479}; \textsc{sha256} prefix \texttt{e7eb512a2ec2438d}.
\par
\textbf{Used in this report by} \hyperref[conj:extcb]{\texttt{conj-\allowbreak{}extcb}}, \hyperref[conj:hcb]{\texttt{conj-\allowbreak{}hcb}}, \hyperref[lem:compcb-amplified-almost-containment]{\texttt{lem-\allowbreak{}compcb-\allowbreak{}amplified-\allowbreak{}almost-\allowbreak{}containment}}, \hyperref[lem:compcb-amplified-compression]{\texttt{lem-\allowbreak{}compcb-\allowbreak{}amplified-\allowbreak{}compression}}, \hyperref[lem:compcb-amplified-compression-identities]{\texttt{lem-\allowbreak{}compcb-\allowbreak{}amplified-\allowbreak{}compression-\allowbreak{}identities}}, \hyperref[lem:compcb-compressed-unit-action]{\texttt{lem-\allowbreak{}compcb-\allowbreak{}compressed-\allowbreak{}unit-\allowbreak{}action}} and 10 further reproduced results (20 registry results import it in total; see \texttt{argument/INDEX.md}).
\par
\textbf{Canonical shard.} \texttt{definitions/\allowbreak{}def-\allowbreak{}extended-\allowbreak{}epsilon-\allowbreak{}cstar-\allowbreak{}algebra.\allowbreak{}md}; twins: \texttt{definitions/INDEX.md}, \texttt{argument/DAG.md}.
\end{defnote}

%% ---- def-fd-cstar-diagonal ---------------------------------------------
\hypertarget{def:fd-cstar-diagonal}{}%
\begin{definition}[diagonal of a finite-dimensional $C^{*}$-algebra\normalfont{} (\texttt{def-\allowbreak{}fd-\allowbreak{}cstar-\allowbreak{}diagonal})]\label{def:fd-cstar-diagonal}
\emph{Aliases:} finite-dimensional $C^{*}$-algebra diagonal; algebra diagonal\par\smallskip
\textbf{Byte-verbatim source text.}

\begin{defsourcequote}
\defsourceline{For arbitrary $\tilde{D}=\sum_{j}A_j\otimes B_j\in\calB\otimes\calB$ and $X\in\calB$, we write}
\defsourceline{\begin{equation}}
\defsourceline{X\tilde{D}=\sum_{j}XA_j\otimes B_j,\qquad}
\defsourceline{\tilde{D}X=\sum_{j}A_j\otimes B_jX,\qquad}
\defsourceline{\pi(\tilde{D})=\sum_{j}A_jB_j.}
\defsourceline{\end{equation}}
\defsourceline{These operations extend from $\calB\otimes\calB$ to $\calB\hotimes\calB$. An element $D\in\calB\hotimes\calB$ is called a \emph{diagonal} if}
\defsourceline{\begin{equation}}
\defsourceline{XD=DX\quad\text{for all }\, X\in\calB,\qquad \pi(D)=I_\calB.}
\defsourceline{\end{equation}}
\end{defsourcequote}

\textbf{Specialization.} The registry uses this definition only when
\texttt{\textbackslash{}(\textbackslash{}mathcal B\textbackslash{})} is a finite-dimensional \texttt{\textbackslash{}(C\textasciicircum{}*\textbackslash{})}-algebra.

\textbf{Provenance.} Byte-verbatim from the pinned source at the recorded locus and
SHA256 prefix.  User-ratified and locked 2026-07-24 (tobiasosborne, in-session sign-off).
\end{definition}
\begin{defnote}
\textbf{Kind.} cited (byte-matched to a pinned source under \texttt{refs/}); \texttt{status: locked}.
\par
\textbf{Provenance.} \texttt{source: kitaev-\allowbreak{}2405.\allowbreak{}02434}; locus \texttt{approximate\_algebras.\allowbreak{}tex:\allowbreak{}1233-\allowbreak{}1242}; \textsc{sha256} prefix \texttt{e7eb512a2ec2438d}.
\par
\textbf{Used by} 3 registry results, none of them reproduced in this report (exploration track; see \texttt{argument/INDEX.md} and \texttt{report/UNWIRED.md}).
\par
\textbf{Canonical shard.} \texttt{definitions/\allowbreak{}def-\allowbreak{}fd-\allowbreak{}cstar-\allowbreak{}diagonal.\allowbreak{}md}; twins: \texttt{definitions/INDEX.md}, \texttt{argument/DAG.md}.
\end{defnote}

%% ---- def-h-space-left-inversion ----------------------------------------
\hypertarget{def:h-space-left-inversion}{}%
\begin{definition}[H-space and left inversion\normalfont{} (\texttt{def-\allowbreak{}h-\allowbreak{}space-\allowbreak{}left-\allowbreak{}inversion})]\label{def:h-space-left-inversion}
\emph{Aliases:} H-space; associative H-space; inversion map; left inversion map\par\smallskip
\textbf{Byte-verbatim source text.}

\begin{defsourcequote}
\defsourceline{By definition, an \emph{H-space} is a topological space $M$ with a basepoint $e$ and a continuous multiplication map $\mu\colon M\to M$ such that $\mu(e,e)=e$ and}
\defsourceline{\begin{equation}\label{H-unit}}
\defsourceline{\bigl(x\mapsto \mu(x,e)\bigr)}
\defsourceline{\,\sim\,\bigl(x\mapsto x\bigr)}
\defsourceline{\,\sim\,\bigl(x\mapsto \mu(e,x)\bigr),}
\defsourceline{\end{equation}}
\defsourceline{where ``$f\sim g$'' means that $f$ and $g$ are homotopic through basepoint-preserving maps. [...] An H-space $M$ is called \emph{associative} if}
\defsourceline{\begin{equation}}
\defsourceline{\bigl((x,y,z)\mapsto \mu(\mu(x,y),z)\bigr)}
\defsourceline{\,\sim\,\bigl((x,y,z)\mapsto \mu(x,\mu(y,z))\bigr).}
\defsourceline{\end{equation}}
\defsourceline{[...] An \emph{inversion map} is a continuous basepoint-preserving map $\sigma\colon M\to M$ such that}
\defsourceline{\begin{equation}\label{homo_inverse}}
\defsourceline{\bigl(x\mapsto \mu(\sigma(x),x)\bigr)}
\defsourceline{\,\sim\,\bigl(x\mapsto e\bigr)}
\defsourceline{\,\sim\,\bigl(x\mapsto \mu(x,\sigma(x))\bigr).}
\defsourceline{\end{equation}}
\defsourceline{[...] When only the existence of that homotopy is required, $\sigma$ is called a \emph{left inversion map}.}
\end{defsourcequote}

\textbf{Statement (harmonised).} An \emph{H-space} is a topological space $M$ with basepoint
$e$ and a continuous multiplication $\mu\colon M\times M\to M$ with $\mu(e,e)=e$
such that $x\mapsto\mu(x,e)$ and $x\mapsto\mu(e,x)$ are each basepoint-preserving
homotopic to the identity. It is \emph{associative} if the two triple-product maps are
homotopic. An \emph{inversion map} is a continuous basepoint-preserving
$\sigma\colon M\to M$ with both $x\mapsto\mu(\sigma(x),x)$ and
$x\mapsto\mu(x,\sigma(x))$ homotopic to the constant map $e$; when only the first
homotopy is required, $\sigma$ is a \emph{left inversion map}.

\textbf{Notes / provenance.} Definition only (per the DESIGN-FUDW-DECOMP-v4.1 \S{}4.1
register row: ``pinned TeX 895--912, definition only'') --- no claim about $\mathcal U$
or any Stage-1 object is carried by this shard. Source typo note: the source
writes $\mu\colon M\to M$ for the multiplication map's type where
$\mu\colon M\times M\to M$ is meant (its uses $\mu(x,y)$ are two-argument
throughout); the harmonised statement records the corrected type. Elisions
\texttt{[...\allowbreak{}]} in the verbatim block skip source sentences about the specific space
$\mathcal U$ (claims, not definition content). Used by the Stage-1 topology rows
\texttt{lem-\allowbreak{}topology-\allowbreak{}hopf-\allowbreak{}structure} and (via fixed-point data) the Lefschetz rows; see
\hyperref[def:lefschetz-fixed-point-data]{\texttt{def-\allowbreak{}lefschetz-\allowbreak{}fixed-\allowbreak{}point-\allowbreak{}data}}.
\end{definition}
\begin{defnote}
\textbf{Kind.} cited (byte-matched to a pinned source under \texttt{refs/}); \texttt{status: locked}.
\par
\textbf{Provenance.} \texttt{source: kitaev-\allowbreak{}2405.\allowbreak{}02434}; locus \texttt{approximate\_algebras.\allowbreak{}tex:\allowbreak{}895-\allowbreak{}912}; \textsc{sha256} prefix \texttt{e7eb512a2ec2438d}; record: byte-\allowbreak{}match lock per definitions/\allowbreak{}README.\allowbreak{}md (cited criterion);\allowbreak{} provisioned by DESIGN-\allowbreak{}FUDW-\allowbreak{}DECOMP-\allowbreak{}v4.\allowbreak{}1.\allowbreak{}md \S{}4.\allowbreak{}1 (kind "cited /\allowbreak{} lock after byte-\allowbreak{}match").
\par
\textbf{Used by} 1 registry result, none of them reproduced in this report (exploration track; see \texttt{argument/INDEX.md} and \texttt{report/UNWIRED.md}).
\par
\textbf{Canonical shard.} \texttt{definitions/\allowbreak{}def-\allowbreak{}h-\allowbreak{}space-\allowbreak{}left-\allowbreak{}inversion.\allowbreak{}md}; twins: \texttt{definitions/INDEX.md}, \texttt{argument/DAG.md}.
\end{defnote}

%% ---- def-ha-map --------------------------------------------------------
\hypertarget{def:ha-map}{}%
\begin{definition}[Ha-map\normalfont{} (\texttt{def-\allowbreak{}ha-\allowbreak{}map})]\label{def:ha-map}
\emph{Aliases:} Ha map; Kitaev Ha-map\par\smallskip
\textbf{Byte-verbatim source text.}

\begin{defsourcequote}
\defsourceline{Let $Q$ be a one-dimensional $\delta$-projection in an $\eps$-$C^*$ algebra $\calA$. For any $\delta$-projections $P,R\in\calA$, we define a map $\Ha^{Q}_{P,R}\colon\calS_{P,R}\to\Bo(\calS_{R,Q},\calS_{P,Q})$. (Here $\Bo(\calS_{R,Q},\calS_{P,Q})$ is the space of bounded linear maps from $\calS_{R,Q}$ to $\calS_{P,Q}$ with the norm induced by the Euclidean norm $\|X\|_\Euc=\sqrt{\braket{X}{X}}$. Due to inequality \eqref{1dQ_ip_norm}, the latter is equal to $\|X\|$ up to a $1\pm O(\eps+\delta)$ factor.) For each $Z\in\calS_{P,R}$ and $X\in\calS_{R,Q}$, the element $\Ha^{Q}_{P,R}(Z)(X)\in\calS_{P,Q}$ is defined by the condition}
\defsourceline{\begin{equation}\label{Ha_def}}
\defsourceline{(Y^\dag\cdot Z)\cdot X+Y^\dag\cdot(Z\cdot X)}
\defsourceline{=2\,\bbraket{Y}{\Ha^{Q}_{P,R}(Z)(X)}\,\wt{Q}\quad \text{for all }\,Y\in\calS_{P,Q}.}
\defsourceline{\end{equation}}
\end{defsourcequote}

\textbf{Notation.} Registry contracts use \texttt{Ha} for the source macro \texttt{\textbackslash{}(\textbackslash{}Ha\textbackslash{})}.

\textbf{Provenance.} Byte-verbatim from the pinned source at the recorded locus and
SHA256 prefix.  User-ratified and locked 2026-07-24 (tobiasosborne, in-session sign-off).
\end{definition}
\begin{defnote}
\textbf{Kind.} cited (byte-matched to a pinned source under \texttt{refs/}); \texttt{status: locked}.
\par
\textbf{Provenance.} \texttt{source: kitaev-\allowbreak{}2405.\allowbreak{}02434}; locus \texttt{approximate\_algebras.\allowbreak{}tex:\allowbreak{}1146-\allowbreak{}1149}; \textsc{sha256} prefix \texttt{e7eb512a2ec2438d}.
\par
\textbf{Used in this report by} \hyperref[conj:extcb]{\texttt{conj-\allowbreak{}extcb}}, \hyperref[conj:hcb]{\texttt{conj-\allowbreak{}hcb}}, \hyperref[lem:hcb1-column-action]{\texttt{lem-\allowbreak{}hcb1-\allowbreak{}column-\allowbreak{}action}}, \hyperref[lem:hcb1-variational-identity]{\texttt{lem-\allowbreak{}hcb1-\allowbreak{}variational-\allowbreak{}identity}}, \hyperref[lem:hcb2-amplified-adjointness]{\texttt{lem-\allowbreak{}hcb2-\allowbreak{}amplified-\allowbreak{}adjointness}}, \hyperref[lem:hcb2-product-defect]{\texttt{lem-\allowbreak{}hcb2-\allowbreak{}product-\allowbreak{}defect}} and 7 further reproduced results (13 registry results import it in total; see \texttt{argument/INDEX.md}).
\par
\textbf{Canonical shard.} \texttt{definitions/\allowbreak{}def-\allowbreak{}ha-\allowbreak{}map.\allowbreak{}md}; twins: \texttt{definitions/INDEX.md}, \texttt{argument/DAG.md}.
\end{defnote}

%% ---- def-lefschetz-fixed-point-data ------------------------------------
\hypertarget{def:lefschetz-fixed-point-data}{}%
\begin{definition}[Lefschetz number and local fixed-point index\normalfont{} (\texttt{def-\allowbreak{}lefschetz-\allowbreak{}fixed-\allowbreak{}point-\allowbreak{}data})]\label{def:lefschetz-fixed-point-data}
\emph{Aliases:} Lefschetz number; fixed point index; fixed-point index; ind(f,x); Lambda(f)\par\smallskip
\textbf{Byte-verbatim source text.}

\begin{defsourcequote}
\defsourceline{The \emph{fixed point index} $\ind(f,x)$ and the \emph{Lefschetz number} $\Lambda(f)$ are defined as follows:}
\defsourceline{\begin{alignat}{2}}
\defsourceline{\ind(f,x)&=\sgn\det(1-f'(x)),\qquad}
\defsourceline{& f'(x)\colon \Ta_xM&\to\Ta_xM,\\[2pt]}
\defsourceline{\Lambda(f)&=\sum_{k}(-1)^k\Tr f^{*k},\qquad}
\defsourceline{& f^{*k}\colon H^k(M;\RR)&\to H^k(M;\RR),}
\defsourceline{\end{alignat}}
\end{defsourcequote}

\textbf{Statement (harmonised).} For a smooth self-map $f$ of a compact orientable
manifold $M$: at a fixed point $x$ with $\det(1-f'(x))\ne0$ (where
$f'(x)\colon T_xM\to T_xM$ is the differential), the \emph{fixed point index} is
$\operatorname{ind}(f,x)=\operatorname{sgn}\det(1-f'(x))$; the \emph{Lefschetz number}
is $\Lambda(f)=\sum_k(-1)^k\operatorname{Tr}f^{*k}$, where
$f^{*k}\colon H^k(M;\mathbb R)\to H^k(M;\mathbb R)$ is the induced map on real
cohomology.

\textbf{Notes / provenance.} Definitions only (per the DESIGN-FUDW-DECOMP-v4.1 \S{}4.1
register row: ``pinned TeX 957--967, definitions only''): this shard carries the two
displayed definitions and NO theorem content --- in particular it does not assert
the Lefschetz--Hopf identity $\sum_{x\in\operatorname{Fix}(f)}\operatorname{ind}(f,x)=\Lambda(f)$
(that is the separate result row \texttt{lem-\allowbreak{}topology-\allowbreak{}lefschetz-\allowbreak{}hopf}, to be \texttt{cited}
from Arkowitz--Brown), nor the index-sign theorem (\texttt{lem-\allowbreak{}topology-\allowbreak{}local-\allowbreak{}index-\allowbreak{}sign},
Granas--Dugundji). The source's index formula is stated in the nondegenerate
($\det\ne0$) smooth case; the general topological index the theorem rows use is
the standard one in their own sources. See \hyperref[def:h-space-left-inversion]{\texttt{def-\allowbreak{}h-\allowbreak{}space-\allowbreak{}left-\allowbreak{}inversion}}.
\end{definition}
\begin{defnote}
\textbf{Kind.} cited (byte-matched to a pinned source under \texttt{refs/}); \texttt{status: locked}.
\par
\textbf{Provenance.} \texttt{source: kitaev-\allowbreak{}2405.\allowbreak{}02434}; locus \texttt{approximate\_algebras.\allowbreak{}tex:\allowbreak{}961-\allowbreak{}967}; \textsc{sha256} prefix \texttt{e7eb512a2ec2438d}; record: byte-\allowbreak{}match lock per definitions/\allowbreak{}README.\allowbreak{}md (cited criterion);\allowbreak{} provisioned by DESIGN-\allowbreak{}FUDW-\allowbreak{}DECOMP-\allowbreak{}v4.\allowbreak{}1.\allowbreak{}md \S{}4.\allowbreak{}1 (kind "cited /\allowbreak{} lock after byte-\allowbreak{}match").
\par
\textbf{Used by} 2 registry results, none of them reproduced in this report (exploration track; see \texttt{argument/INDEX.md} and \texttt{report/UNWIRED.md}).
\par
\textbf{Canonical shard.} \texttt{definitions/\allowbreak{}def-\allowbreak{}lefschetz-\allowbreak{}fixed-\allowbreak{}point-\allowbreak{}data.\allowbreak{}md}; twins: \texttt{definitions/INDEX.md}, \texttt{argument/DAG.md}.
\end{defnote}

%% ---- def-one-dimensional-delta-projection ------------------------------
\hypertarget{def:one-dimensional-delta-projection}{}%
\begin{definition}[one-dimensional $\delta$-projection and equivalence\normalfont{} (\texttt{def-\allowbreak{}one-\allowbreak{}dimensional-\allowbreak{}delta-\allowbreak{}projection})]\label{def:one-dimensional-delta-projection}
\emph{Aliases:} one-dimensional delta-projection; equivalent delta-projections\par\smallskip
\textbf{Byte-verbatim source text.}

\begin{defsourcequote}
\defsourceline{When $P=Q$, the abbreviations $\Co_{P}$, $\calS_P$ are used. It is clear that $\calS_P=0$ if and only if $P$ is sufficiently close to $0$ as described by the first alternative in \eqref{P_alternatives}. (Recall that in the opposite case, we call $P$ ``nonvanishing''.) Similarly, $\Co_P=1$ and $\calS_P=\calA$ if and only if $P$ is close to $I$. If $\dim\calS_P=1$, we say that $P$ is a \emph{one-dimensional $\delta$-projection}.}
\end{defsourcequote}

\begin{defsourcequote}
\defsourceline{One-dimensional $\delta$-projections $P$ and $Q$ are called \emph{equivalent} if $\dim\calS_{P,Q}=1$. This is indeed an equivalence relation; in particular, it is transitive due to Lemma~\ref{lem_PQR}.}
\end{defsourcequote}

\textbf{Status.} Locked after the cited-definition byte-match criterion was user-ratified on 2026-07-24.

\textbf{Provenance.} Byte-verbatim from the pinned source at the recorded loci
and SHA256 prefix; provisioned by \texttt{DESIGN-\allowbreak{}FUDW-\allowbreak{}DECOMP-\allowbreak{}v3.\allowbreak{}md} \S{}4.1 and admitted
only for the safe subset by \texttt{VERDICT-\allowbreak{}FUDW-\allowbreak{}DECOMP-\allowbreak{}V3.\allowbreak{}md} \S{}D.
\end{definition}
\begin{defnote}
\textbf{Kind.} cited (byte-matched to a pinned source under \texttt{refs/}); \texttt{status: locked}.
\par
\textbf{Provenance.} \texttt{source: kitaev-\allowbreak{}2405.\allowbreak{}02434}; locus \texttt{approximate\_algebras.\allowbreak{}tex:\allowbreak{}1066,\allowbreak{}1187}; \textsc{sha256} prefix \texttt{e7eb512a2ec2438d}; record: user-\allowbreak{}ratified 2026-\allowbreak{}07-\allowbreak{}24 (tobiasosborne,\allowbreak{} in-\allowbreak{}session sign-\allowbreak{}off:\allowbreak{} byte-\allowbreak{}match-\allowbreak{}to-\allowbreak{}source criterion for cited;\allowbreak{} delegated ratification for consensus/\allowbreak{}original).
\par
\textbf{Used in this report by} \hyperref[conj:hcb]{\texttt{conj-\allowbreak{}hcb}}, \hyperref[lem:extcb-one-dimensional-corner-dimension]{\texttt{lem-\allowbreak{}extcb-\allowbreak{}one-\allowbreak{}dimensional-\allowbreak{}corner-\allowbreak{}dimension}}, \hyperref[lem:extcb-one-dimensional-product]{\texttt{lem-\allowbreak{}extcb-\allowbreak{}one-\allowbreak{}dimensional-\allowbreak{}product}} (3 registry results import it in total; see \texttt{argument/INDEX.md}).
\par
\textbf{Canonical shard.} \texttt{definitions/\allowbreak{}def-\allowbreak{}one-\allowbreak{}dimensional-\allowbreak{}delta-\allowbreak{}projection.\allowbreak{}md}; twins: \texttt{definitions/INDEX.md}, \texttt{argument/DAG.md}.
\end{defnote}

%% ---- def-projection-basis ----------------------------------------------
\hypertarget{def:projection-basis}{}%
\begin{definition}[projection basis of a finite-dimensional commutative $C^{*}$-algebra\normalfont{} (\texttt{def-\allowbreak{}projection-\allowbreak{}basis})]\label{def:projection-basis}
\emph{Aliases:} projection basis\par\smallskip
\textbf{Byte-verbatim source text.}

\begin{defsourcequote}
\defsourceline{The second application of Lemma~\ref{lem_merging} involves an auxiliary statement about $\delta$-projections $P_1,\dots,P_p,Q_1,\dots,Q_q$ satisfying certain conditions, which will be formulated in terms of inclusions of finite-dimensional commutative $C^*$ algebras. Such an algebra $\calB$ is described by a \emph{projection basis} $\{\Pi_1,\dots,\Pi_n\}\subseteq\calB$ such that $\Pi_j^\dag=\Pi_j$,\, $\Pi_j\Pi_k=\delta_{jk}\Pi_k$, and $\sum_{j=1}^{n}\Pi_j=I$.}
\end{defsourcequote}

\textbf{Status.} Draft transcription; ratification is required before locking.

\textbf{Provenance.} Byte-verbatim from the pinned source at the recorded locus
and SHA256 prefix; provisioned by \texttt{DESIGN-\allowbreak{}FUDW-\allowbreak{}DECOMP-\allowbreak{}v3.\allowbreak{}md} \S{}4.1 and admitted
only for the safe subset by \texttt{VERDICT-\allowbreak{}FUDW-\allowbreak{}DECOMP-\allowbreak{}V3.\allowbreak{}md} \S{}D.
\end{definition}
\begin{defnote}
\textbf{Kind.} cited (byte-matched to a pinned source under \texttt{refs/}); \texttt{status: locked}.
\par
\textbf{Provenance.} \texttt{source: kitaev-\allowbreak{}2405.\allowbreak{}02434}; locus \texttt{approximate\_algebras.\allowbreak{}tex:\allowbreak{}1361}; \textsc{sha256} prefix \texttt{e7eb512a2ec2438d}; record: user-\allowbreak{}ratified 2026-\allowbreak{}07-\allowbreak{}24 (tobiasosborne,\allowbreak{} in-\allowbreak{}session sign-\allowbreak{}off:\allowbreak{} byte-\allowbreak{}match-\allowbreak{}to-\allowbreak{}source criterion for cited;\allowbreak{} delegated ratification for consensus/\allowbreak{}original).
\par
\textbf{Used in this report by} \hyperref[lem:extcb-corner-dimension-additivity]{\texttt{lem-\allowbreak{}extcb-\allowbreak{}corner-\allowbreak{}dimension-\allowbreak{}additivity}}, \hyperref[lem:extcb1-cross-corner-dimension]{\texttt{lem-\allowbreak{}extcb1-\allowbreak{}cross-\allowbreak{}corner-\allowbreak{}dimension}} (2 registry results import it in total; see \texttt{argument/INDEX.md}).
\par
\textbf{Canonical shard.} \texttt{definitions/\allowbreak{}def-\allowbreak{}projection-\allowbreak{}basis.\allowbreak{}md}; twins: \texttt{definitions/INDEX.md}, \texttt{argument/DAG.md}.
\end{defnote}

%% ---- def-theta-idempotent-approximation --------------------------------
\hypertarget{def:theta-idempotent-approximation}{}%
\begin{definition}[theta idempotent-approximation map\normalfont{} (\texttt{def-\allowbreak{}theta-\allowbreak{}idempotent-\allowbreak{}approximation})]\label{def:theta-idempotent-approximation}
\emph{Aliases:} theta map; functional-calculus idempotent approximation\par\smallskip
\textbf{Byte-verbatim source text.}

Power-series functional calculus (tex:505-510):

\begin{defsourcequote}
\defsourceline{f(X)=a_0I+\sum_{n=1}^{\infty}a_n(X-x_0I)^n\quad\:}
\defsourceline{\text{if }\, \|X-x_0I\|<r,\\}
\defsourceline{\|f(X)-f(x_0I)\|\le \sum_{n=1}^{\infty}|a_n|\ts\|X-x_0I\|^n.}
\end{defsourcequote}

Absolute value and sign (tex:517-520):

\begin{defsourcequote}
\defsourceline{|X|=(X^2)^{1/2},\qquad \sgn(X)=X(X^2)^{-1/2}\qquad\quad}
\defsourceline{\bigl(\|X^2-x_0^2I\|<x_0^2\bigr).}
\end{defsourcequote}

The theta map (tex:527-528, the definitional display of Proposition \texttt{prop\_P}):

\begin{defsourcequote}
\defsourceline{\wt{P}=\theta(2P-I),\qquad}
\defsourceline{\text{where}\quad \theta(X)=\frac{1}{2}\bigl(I+\sgn(X)\bigr).}
\end{defsourcequote}

\textbf{Scope note.} This is the map \texttt{theta} referenced by the compression-map
display \texttt{Co\_\{P,\allowbreak{}Q\}=\textbackslash{}theta(\textbackslash{}La\_P\textbackslash{}Ra\_Q+\textbackslash{}Ra\_Q\textbackslash{}La\_P-\allowbreak{}1)} in
\hyperref[def:compressed-corner]{\texttt{def-\allowbreak{}compressed-\allowbreak{}corner}} (tex:1057). The neighbouring Proposition's
conclusions (idempotence of \texttt{wtP}, commutation, distance bound) are claims,
not definitions --- they enter the registry only as re-proved results.

\textbf{Ratification.} Draft pending recorded sign-off; expected to lock under the
user's recorded byte-match criterion (2026-07-24: ``if they match the source
then they are valid'').
\end{definition}
\begin{defnote}
\textbf{Kind.} cited (byte-matched to a pinned source under \texttt{refs/}); \texttt{status: draft}.
\par
\textbf{Provenance.} \texttt{source: kitaev-\allowbreak{}2405.\allowbreak{}02434}; locus \texttt{approximate\_algebras.\allowbreak{}tex:\allowbreak{}505-\allowbreak{}528 (displays Taylor\_simple,\allowbreak{} Taylor\_simple\_bound,\allowbreak{} abs\_sgn,\allowbreak{} and the theta display inside Proposition prop\_P ---\allowbreak{} ONLY the definitional displays are cited;\allowbreak{} the Proposition's claims are NOT citable and must be re-\allowbreak{}proved)}; \textsc{sha256} prefix \texttt{e7eb512a2ec2438d}.
\par
\textbf{Used in this report by} \hyperref[lem:compcb-amplification-naturality]{\texttt{lem-\allowbreak{}compcb-\allowbreak{}amplification-\allowbreak{}naturality}}, \hyperref[lem:compcb-entrywise-compression-naturality]{\texttt{lem-\allowbreak{}compcb-\allowbreak{}entrywise-\allowbreak{}compression-\allowbreak{}naturality}} (2 registry results import it in total; see \texttt{argument/INDEX.md}).
\par
\textbf{Canonical shard.} \texttt{definitions/\allowbreak{}def-\allowbreak{}theta-\allowbreak{}idempotent-\allowbreak{}approximation.\allowbreak{}md}; twins: \texttt{definitions/INDEX.md}, \texttt{argument/DAG.md}.
\end{defnote}


% Project-internal packaging: hypothesis data and derived notation
% GENERATED by scripts/gen-report-defs.py — DO NOT HAND-EDIT.
% Single source of truth: definitions/<id>.md (CLAUDE.md L2).  Re-render with
%   python3 scripts/gen-report-defs.py
% The gate `python3 scripts/gen-report-defs.py --check` fails when this file is stale.

% Layer file: Project-internal packaging: hypothesis data and derived notation
\subsection*{Project-internal packaging: hypothesis data and derived notation}

%% ---- def-canonical-corner-identifications ------------------------------
\hypertarget{def:canonical-corner-identifications}{}%
\begin{definition}[canonical corner identifications\normalfont{} (\texttt{def-\allowbreak{}canonical-\allowbreak{}corner-\allowbreak{}identifications})]\label{def:canonical-corner-identifications}
After identifying the one-dimensional corner
\(\mathcal S_{Q,Q}\) with \(\mathbb C\) by the normalized vector \(q_0\), the
canonical column identification is
\[
[J_{P,Q,n}(Z)c]_i=\sum_j Z_{ij}c_j,
\]
and the canonical row identification is
\[
J_{Q,P,n}(Z)=J_{P,Q,n}(Z^\dagger)^\dagger.
\]
\end{definition}
\begin{defnote}
\textbf{Kind.} original (introduced by this project); \texttt{status: locked}.\quad\textbf{Aliases.} canonical J maps; canonical column and row identifications.
\par
\textbf{Provenance.} \texttt{source: internal}; locus \texttt{PROOF-\allowbreak{}W74F-\allowbreak{}E-\allowbreak{}HCB.\allowbreak{}md \S{}8.\allowbreak{}1;\allowbreak{} DESIGN-\allowbreak{}FUDW-\allowbreak{}DECOMP-\allowbreak{}v3.\allowbreak{}md \S{}4.\allowbreak{}1}; no source hash (not a literature transcription); record: user-\allowbreak{}ratified 2026-\allowbreak{}07-\allowbreak{}24 (tobiasosborne,\allowbreak{} in-\allowbreak{}session sign-\allowbreak{}off;\allowbreak{} delegated ratification for consensus/\allowbreak{}original tiers).
\par
\textbf{Used in this report by} \hyperref[conj:hcb]{\texttt{conj-\allowbreak{}hcb}}, \hyperref[lem:hcb4-canonical-closeness]{\texttt{lem-\allowbreak{}hcb4-\allowbreak{}canonical-\allowbreak{}closeness}}, \hyperref[lem:hcb4-canonical-gram]{\texttt{lem-\allowbreak{}hcb4-\allowbreak{}canonical-\allowbreak{}gram}}, \hyperref[lem:hcb4-canonical-inverse]{\texttt{lem-\allowbreak{}hcb4-\allowbreak{}canonical-\allowbreak{}inverse}} (4 registry results import it in total; see \texttt{argument/INDEX.md}).
\par
\textbf{Canonical shard.} \texttt{definitions/\allowbreak{}def-\allowbreak{}canonical-\allowbreak{}corner-\allowbreak{}identifications.\allowbreak{}md}; twins: \texttt{definitions/INDEX.md}, \texttt{argument/DAG.md}.
\par
\textbf{Status.} Draft packaging of notation only; no norm or inverse estimate is
included, and ratification is required before locking.

\textbf{Provenance.} Transcribed from \texttt{PROOF-\allowbreak{}W74F-\allowbreak{}E-\allowbreak{}HCB.\allowbreak{}md} \S{}8.1 and the
provisioning row of \texttt{DESIGN-\allowbreak{}FUDW-\allowbreak{}DECOMP-\allowbreak{}v3.\allowbreak{}md} \S{}4.1; admitted by
\texttt{VERDICT-\allowbreak{}FUDW-\allowbreak{}DECOMP-\allowbreak{}V3.\allowbreak{}md} \S{}D.
\end{defnote}

%% ---- def-compressed-associator -----------------------------------------
\hypertarget{def:compressed-associator}{}%
\begin{definition}[compressed associator\normalfont{} (\texttt{def-\allowbreak{}compressed-\allowbreak{}associator})]\label{def:compressed-associator}
For compatible compressed rectangular products, the
\emph{compressed associator} is
\[
(A\mathbin{\cdot}B)\mathbin{\cdot}C
-A\mathbin{\cdot}(B\mathbin{\cdot}C).
\]
\end{definition}
\begin{defnote}
\textbf{Kind.} original (introduced by this project); \texttt{status: locked}.\quad\textbf{Aliases.} rectangular compressed associator.
\par
\textbf{Provenance.} \texttt{source: internal}; locus \texttt{PROOF-\allowbreak{}W74F-\allowbreak{}E-\allowbreak{}HCB.\allowbreak{}md \S{}3;\allowbreak{} DESIGN-\allowbreak{}FUDW-\allowbreak{}DECOMP-\allowbreak{}v3.\allowbreak{}md \S{}4.\allowbreak{}1}; no source hash (not a literature transcription); record: user-\allowbreak{}ratified 2026-\allowbreak{}07-\allowbreak{}24 (tobiasosborne,\allowbreak{} in-\allowbreak{}session sign-\allowbreak{}off;\allowbreak{} delegated ratification for consensus/\allowbreak{}original tiers).
\par
\textbf{Used in this report by} \hyperref[lem:hcb0-compressed-associator]{\texttt{lem-\allowbreak{}hcb0-\allowbreak{}compressed-\allowbreak{}associator}}, \hyperref[lem:hcb1-column-action]{\texttt{lem-\allowbreak{}hcb1-\allowbreak{}column-\allowbreak{}action}} (2 registry results import it in total; see \texttt{argument/INDEX.md}).
\par
\textbf{Canonical shard.} \texttt{definitions/\allowbreak{}def-\allowbreak{}compressed-\allowbreak{}associator.\allowbreak{}md}; twins: \texttt{definitions/INDEX.md}, \texttt{argument/DAG.md}.
\par
\textbf{Status.} Draft project name for the displayed source expression; it
asserts no estimate and requires ratification before locking.

\textbf{Provenance.} Transcribed from \texttt{PROOF-\allowbreak{}W74F-\allowbreak{}E-\allowbreak{}HCB.\allowbreak{}md} \S{}3 and the provisioning
row of \texttt{DESIGN-\allowbreak{}FUDW-\allowbreak{}DECOMP-\allowbreak{}v3.\allowbreak{}md} \S{}4.1; admitted by
\texttt{VERDICT-\allowbreak{}FUDW-\allowbreak{}DECOMP-\allowbreak{}V3.\allowbreak{}md} \S{}D.
\end{defnote}

%% ---- def-extcb-datum ---------------------------------------------------
\hypertarget{def:extcb-datum}{}%
\begin{definition}[closed EXT-CB datum\normalfont{} (\texttt{def-\allowbreak{}extcb-\allowbreak{}datum})]\label{def:extcb-datum}
A \emph{closed EXT-CB datum} consists of an extended
\(\varepsilon\)-\(C^*\)-algebra \(\mathcal A\); \(\delta\)-projections
\(P,Q\) with \(\lVert P+Q-I\rVert\le\delta\); an extended
\(\delta\)-isomorphism \(v:M_r\to S_P\); \(\dim S_Q=1\); and
\(S_{P,Q}\ne0\), with \(e=\delta+\varepsilon\).
\end{definition}
\begin{defnote}
\textbf{Kind.} original (introduced by this project); \texttt{status: locked}.\quad\textbf{Aliases.} EXT-CB datum.
\par
\textbf{Provenance.} \texttt{source: internal}; locus \texttt{DESIGN-\allowbreak{}FUDW-\allowbreak{}DECOMP-\allowbreak{}v3.\allowbreak{}md \S{}4.\allowbreak{}1 exact proposed datum package}; no source hash (not a literature transcription); record: user-\allowbreak{}ratified 2026-\allowbreak{}07-\allowbreak{}24 (tobiasosborne,\allowbreak{} in-\allowbreak{}session sign-\allowbreak{}off;\allowbreak{} delegated ratification for consensus/\allowbreak{}original tiers).
\par
\textbf{Used in this report by} \hyperref[conj:extcb]{\texttt{conj-\allowbreak{}extcb}}, \hyperref[lem:extcb1-close-corner-dimension]{\texttt{lem-\allowbreak{}extcb1-\allowbreak{}close-\allowbreak{}corner-\allowbreak{}dimension}}, \hyperref[lem:extcb1-cross-corner-dimension]{\texttt{lem-\allowbreak{}extcb1-\allowbreak{}cross-\allowbreak{}corner-\allowbreak{}dimension}} (3 registry results import it in total; see \texttt{argument/INDEX.md}).
\par
\textbf{Canonical shard.} \texttt{definitions/\allowbreak{}def-\allowbreak{}extcb-\allowbreak{}datum.\allowbreak{}md}; twins: \texttt{definitions/INDEX.md}, \texttt{argument/DAG.md}.
\par
\textbf{Status.} Draft theorem-free hypothesis package; it includes no extension
conclusion and requires ratification before locking.

\textbf{Provenance.} Byte-faithful transcription of the exact datum package in
\texttt{DESIGN-\allowbreak{}FUDW-\allowbreak{}DECOMP-\allowbreak{}v3.\allowbreak{}md} \S{}4.1, approved as definition provisioning by
\texttt{VERDICT-\allowbreak{}FUDW-\allowbreak{}DECOMP-\allowbreak{}V3.\allowbreak{}md} \S{}\S{}6,D.
\end{defnote}

%% ---- def-extended-delta-inclusion --------------------------------------
\hypertarget{def:extended-delta-inclusion}{}%
\begin{definition}[extended $\delta$-inclusion and isomorphism\normalfont{} (\texttt{def-\allowbreak{}extended-\allowbreak{}delta-\allowbreak{}inclusion})]\label{def:extended-delta-inclusion}
An \emph{extended \(\delta\)-inclusion} is a linear map \(v\) for
which every amplification \(1_{M_n}\otimes v\) is a \(\delta\)-inclusion:
it is a \(\delta\)-homomorphism and satisfies the two-sided
\((1\pm\delta)\) norm bounds. It is an \emph{extended
\(\delta\)-isomorphism} when \(v\) is bijective.
\end{definition}
\begin{defnote}
\textbf{Kind.} consensus (project-internal, agreed; no single external source); \texttt{status: locked}.\quad\textbf{Aliases.} extended delta-inclusion; extended delta-isomorphism.
\par
\textbf{Provenance.} \texttt{source: internal}; locus \texttt{harmonization of approximate\_algebras.\allowbreak{}tex:\allowbreak{}443-\allowbreak{}456,\allowbreak{}1477-\allowbreak{}1484 prescribed by DESIGN-\allowbreak{}FUDW-\allowbreak{}DECOMP-\allowbreak{}v3.\allowbreak{}md \S{}4.\allowbreak{}1}; no source hash (not a literature transcription); record: user-\allowbreak{}ratified 2026-\allowbreak{}07-\allowbreak{}24 (tobiasosborne,\allowbreak{} in-\allowbreak{}session sign-\allowbreak{}off;\allowbreak{} delegated ratification for consensus/\allowbreak{}original tiers).
\par
\textbf{Used in this report by} \hyperref[conj:extcb]{\texttt{conj-\allowbreak{}extcb}}, \hyperref[lem:compcb-single-compression-transfer]{\texttt{lem-\allowbreak{}compcb-\allowbreak{}single-\allowbreak{}compression-\allowbreak{}transfer}}, \hyperref[lem:extcb-corner-dimension-additivity]{\texttt{lem-\allowbreak{}extcb-\allowbreak{}corner-\allowbreak{}dimension-\allowbreak{}additivity}}, \hyperref[lem:extcb-four-corner-merge]{\texttt{lem-\allowbreak{}extcb-\allowbreak{}four-\allowbreak{}corner-\allowbreak{}merge}} (5 registry results import it in total; see \texttt{argument/INDEX.md}).
\par
\textbf{Canonical shard.} \texttt{definitions/\allowbreak{}def-\allowbreak{}extended-\allowbreak{}delta-\allowbreak{}inclusion.\allowbreak{}md}; twins: \texttt{definitions/INDEX.md}, \texttt{argument/DAG.md}.
\par
\textbf{Status.} Draft harmonization only; it contains no theorem and requires
ratification before locking.

\textbf{Provenance.} Transcribed from the exact provisioning decision in
\texttt{DESIGN-\allowbreak{}FUDW-\allowbreak{}DECOMP-\allowbreak{}v3.\allowbreak{}md} \S{}4.1, which harmonizes the two pinned-source loci;
\texttt{VERDICT-\allowbreak{}FUDW-\allowbreak{}DECOMP-\allowbreak{}V3.\allowbreak{}md} \S{}\S{}6,D admits it only for the safe subset.
\end{defnote}

%% ---- def-four-corner-merging-datum -------------------------------------
\hypertarget{def:four-corner-merging-datum}{}%
\begin{definition}[amplified four-corner merging datum\normalfont{} (\texttt{def-\allowbreak{}four-\allowbreak{}corner-\allowbreak{}merging-\allowbreak{}datum})]\label{def:four-corner-merging-datum}
An \emph{amplified four-corner merging datum} consists of two
complementary source projections, two approximately complementary target
projections, and four fixed level-one corner maps whose amplifications share
one defect bound for: involution compatibility, compressed-product
compatibility, diagonal-unit approximation, and two-sided norm control. The
four maps are the same maps at every amplification. \textbf{Quantitative
complementarity (amendment 2026-07-25, per the cited locus):} the target
projections are delta-projections and satisfy \(\|P_1+P_2-I\|\le\rho\)
with \(\rho\) the same common defect bound --- byte-faithful to the pinned
source at \texttt{approximate\_algebras.\allowbreak{}tex:\allowbreak{}1326} (``let \(P_1,P_2\) be
\(\delta\)-projections in an \(\varepsilon\)-\(C^*\) algebra \(\mathcal{A}\)
such that \(\|P_1+P_2-I\|\le\delta\)'').
\end{definition}
\begin{defnote}
\textbf{Kind.} original (introduced by this project); \texttt{status: locked}.\quad\textbf{Aliases.} four-corner merging datum.
\par
\textbf{Provenance.} \texttt{source: internal}; locus \texttt{project packaging of approximate\_algebras.\allowbreak{}tex:\allowbreak{}1325-\allowbreak{}1345 prescribed by DESIGN-\allowbreak{}FUDW-\allowbreak{}DECOMP-\allowbreak{}v3.\allowbreak{}md \S{}4.\allowbreak{}1}; no source hash (not a literature transcription); record: user-\allowbreak{}ratified 2026-\allowbreak{}07-\allowbreak{}24 (tobiasosborne,\allowbreak{} in-\allowbreak{}session sign-\allowbreak{}off;\allowbreak{} delegated ratification for consensus/\allowbreak{}original tiers);\allowbreak{} quantitative-\allowbreak{}complementarity amendment user-\allowbreak{}ratified 2026-\allowbreak{}07-\allowbreak{}25 (tobiasosborne,\allowbreak{} in-\allowbreak{}session sign-\allowbreak{}off ---\allowbreak{} transcription-\allowbreak{}fidelity correction to the cited locus;\allowbreak{} af challenge ch-\allowbreak{}4dd98c4d).
\par
\textbf{Used in this report by} \hyperref[lem:extcb-four-corner-merge]{\texttt{lem-\allowbreak{}extcb-\allowbreak{}four-\allowbreak{}corner-\allowbreak{}merge}}, \hyperref[lem:extcb-four-corner-norm]{\texttt{lem-\allowbreak{}extcb-\allowbreak{}four-\allowbreak{}corner-\allowbreak{}norm}} (2 registry results import it in total; see \texttt{argument/INDEX.md}).
\par
\textbf{Canonical shard.} \texttt{definitions/\allowbreak{}def-\allowbreak{}four-\allowbreak{}corner-\allowbreak{}merging-\allowbreak{}datum.\allowbreak{}md}; twins: \texttt{definitions/INDEX.md}, \texttt{argument/DAG.md}.
\par
\textbf{Status.} Draft project packaging of the four source hypotheses; it asserts
no merge conclusion and requires ratification before locking.

\textbf{Provenance.} Transcribed from \texttt{DESIGN-\allowbreak{}FUDW-\allowbreak{}DECOMP-\allowbreak{}v3.\allowbreak{}md} \S{}4.1 and the four
conditions at the pinned-source locus; admitted only for the safe EXT front
end by \texttt{VERDICT-\allowbreak{}FUDW-\allowbreak{}DECOMP-\allowbreak{}V3.\allowbreak{}md} \S{}D.
\end{defnote}

%% ---- def-hcb-datum -----------------------------------------------------
\hypertarget{def:hcb-datum}{}%
\begin{definition}[closed H-CB datum\normalfont{} (\texttt{def-\allowbreak{}hcb-\allowbreak{}datum})]\label{def:hcb-datum}
A \emph{closed H-CB datum} consists of an extended
\(\varepsilon\)-\(C^*\)-algebra \(\mathcal A\); one level-one
one-dimensional \(\delta\)-projection \(Q\); \(\delta\)-projections
\(P,R,S\); the corresponding compressed corners, amplified column spaces,
and their defined sesquilinear forms; and \(e=\delta+\varepsilon\).
\end{definition}
\begin{defnote}
\textbf{Kind.} original (introduced by this project); \texttt{status: locked}.\quad\textbf{Aliases.} H-CB datum.
\par
\textbf{Provenance.} \texttt{source: internal}; locus \texttt{DESIGN-\allowbreak{}FUDW-\allowbreak{}DECOMP-\allowbreak{}v3.\allowbreak{}md \S{}4.\allowbreak{}1 exact proposed datum package}; no source hash (not a literature transcription); record: user-\allowbreak{}ratified 2026-\allowbreak{}07-\allowbreak{}24 (tobiasosborne,\allowbreak{} in-\allowbreak{}session sign-\allowbreak{}off;\allowbreak{} delegated ratification for consensus/\allowbreak{}original tiers).
\par
\textbf{Used in this report by} \hyperref[lem:hcb-column-hilbert-squared]{\texttt{lem-\allowbreak{}hcb-\allowbreak{}column-\allowbreak{}hilbert-\allowbreak{}squared}}, \hyperref[lem:hcb0-compressed-associator]{\texttt{lem-\allowbreak{}hcb0-\allowbreak{}compressed-\allowbreak{}associator}}, \hyperref[lem:hcb1-column-action]{\texttt{lem-\allowbreak{}hcb1-\allowbreak{}column-\allowbreak{}action}}, \hyperref[lem:hcb1-variational-identity]{\texttt{lem-\allowbreak{}hcb1-\allowbreak{}variational-\allowbreak{}identity}}, \hyperref[lem:hcb2-amplified-adjointness]{\texttt{lem-\allowbreak{}hcb2-\allowbreak{}amplified-\allowbreak{}adjointness}}, \hyperref[lem:hcb2-product-defect]{\texttt{lem-\allowbreak{}hcb2-\allowbreak{}product-\allowbreak{}defect}} and 9 further reproduced results (15 registry results import it in total; see \texttt{argument/INDEX.md}).
\par
\textbf{Canonical shard.} \texttt{definitions/\allowbreak{}def-\allowbreak{}hcb-\allowbreak{}datum.\allowbreak{}md}; twins: \texttt{definitions/INDEX.md}, \texttt{argument/DAG.md}.
\par
\textbf{Scope.} This is data and notation only. It includes no column-norm
comparison, Ha estimate, inverse, or parent H-CB conclusion.

\textbf{Status.} Draft theorem-free package; ratification is required before
locking.

\textbf{Provenance.} Byte-faithful transcription of the exact datum package in
\texttt{DESIGN-\allowbreak{}FUDW-\allowbreak{}DECOMP-\allowbreak{}v3.\allowbreak{}md} \S{}4.1, approved as definition provisioning by
\texttt{VERDICT-\allowbreak{}FUDW-\allowbreak{}DECOMP-\allowbreak{}V3.\allowbreak{}md} \S{}\S{}6,D.
\end{defnote}

%% ---- def-positive-approximate-retract ----------------------------------
\hypertarget{def:positive-approximate-retract}{}%
\begin{definition}[positive approximate retract\normalfont{} (\texttt{def-\allowbreak{}positive-\allowbreak{}approximate-\allowbreak{}retract})]\label{def:positive-approximate-retract}
A \emph{positive approximate retract} between finite commutative
$C^*$-algebras is a pair of positive unital maps
\[
A:\ell^\infty(k)\longrightarrow\ell^\infty(n),
\qquad
M:\ell^\infty(n)\longrightarrow\ell^\infty(k)
\]
for which $MA$ is close to $I_k$ in the
$\ell^\infty\!\to\!\ell^\infty$ operator norm.  Its retract defect is
\[
\varepsilon_{\mathrm{ret}}:=\lVert MA-I_k\rVert_{\infty\to\infty}.
\]
Equivalently, the matrices of $A$ and $M$ have probability-vector rows.
\end{definition}
\begin{defnote}
\textbf{Kind.} original (introduced by this project); \texttt{status: locked}.\quad\textbf{Aliases.} positive retract pair; PRH datum.
\par
\textbf{Provenance.} \texttt{source: internal}; locus \texttt{project formulation transcribed from docs/\allowbreak{}plans/\allowbreak{}2026-\allowbreak{}07-\allowbreak{}23-\allowbreak{}W74F-\allowbreak{}artifacts/\allowbreak{}PROOF-\allowbreak{}W74F-\allowbreak{}A-\allowbreak{}PRH.\allowbreak{}md \S{}\S{}1,\allowbreak{}4}; no source hash (not a literature transcription); record: user-\allowbreak{}ratified 2026-\allowbreak{}07-\allowbreak{}24 (tobiasosborne,\allowbreak{} in-\allowbreak{}session sign-\allowbreak{}off).
\par
\textbf{Used in this report by} \hyperref[lem:prh]{\texttt{lem-\allowbreak{}prh}}, \hyperref[lem:routef-prh-finish]{\texttt{lem-\allowbreak{}routef-\allowbreak{}prh-\allowbreak{}finish}} (3 registry results import it in total; see \texttt{argument/INDEX.md}).
\par
\textbf{Canonical shard.} \texttt{definitions/\allowbreak{}def-\allowbreak{}positive-\allowbreak{}approximate-\allowbreak{}retract.\allowbreak{}md}; twins: \texttt{definitions/INDEX.md}, \texttt{argument/DAG.md}.
\par
\textbf{Scope.} This project term packages the datum hardened by \texttt{lem-\allowbreak{}prh}; it
does not assert that such a pair exists for any particular almost-idempotent
map.

\textbf{Provenance.} Project formulation extracted from the W74F-A PRH artifact.
User-ratified and locked 2026-07-24 (tobiasosborne, in-session sign-off).
\end{defnote}

%% ---- def-ucp-map -------------------------------------------------------
\hypertarget{def:ucp-map}{}%
\begin{definition}[unital completely positive map\normalfont{} (\texttt{def-\allowbreak{}ucp-\allowbreak{}map})]\label{def:ucp-map}
Let $A$ and $B$ be unital $C^*$-algebras (in this repo's
Route-F usage: finite-dimensional, typically $\ell_\infty^n$, $M_n$, or a
finite-dimensional unital $C^*$-algebra $\mathcal B$). A linear map
$\varphi:A\to B$ is \emph{positive} if $\varphi(a)\ge0$ whenever $a\ge0$;
\emph{completely positive} (CP) if for every $r\ge1$ the amplification
$\mathrm{id}_{M_r}\otimes\varphi: M_r(A)\to M_r(B)$ is positive; and
\emph{unital} if $\varphi(1_A)=1_B$. A \emph{UCP map} is a unital completely positive
map. Standard facts used silently at the BSc/MSc level: a UCP map is
contractive; a positive map out of a commutative $C^*$-algebra is
automatically CP; compositions of UCP maps are UCP.
\end{definition}
\begin{defnote}
\textbf{Kind.} consensus (project-internal, agreed; no single external source); \texttt{status: locked}.\quad\textbf{Aliases.} UCP map; UCP; completely positive unital map.
\par
\textbf{Provenance.} \texttt{source: internal}; locus \texttt{standard operator-\allowbreak{}algebra textbook notion (e.\allowbreak{}g.\allowbreak{} Paulsen,\allowbreak{} Completely Bounded Maps and Operator Algebras,\allowbreak{} ch.\allowbreak{} 2-\allowbreak{}3);\allowbreak{} adopted for the Route-\allowbreak{}F rows per AUDIT-\allowbreak{}F0-\allowbreak{}ASSEMBLY.\allowbreak{}md sect 1.\allowbreak{}1}; no source hash (not a literature transcription); record: user-\allowbreak{}ratified 2026-\allowbreak{}07-\allowbreak{}27 (W79 decision D3,\allowbreak{} docs/\allowbreak{}plans/\allowbreak{}2026-\allowbreak{}07-\allowbreak{}27-\allowbreak{}W78-\allowbreak{}ratification-\allowbreak{}package.\allowbreak{}md;\allowbreak{} def shard chosen over an L2 textbook exemption because the F0 lift row's contract concludes "Phi is UCP").
\par
\textbf{Used by} 3 registry results, none of them reproduced in this report (exploration track; see \texttt{argument/INDEX.md} and \texttt{report/UNWIRED.md}).
\par
\textbf{Canonical shard.} \texttt{definitions/\allowbreak{}def-\allowbreak{}ucp-\allowbreak{}map.\allowbreak{}md}; twins: \texttt{definitions/INDEX.md}, \texttt{argument/DAG.md}.
\par
\textbf{Notes / provenance.} Provisioned per the F0 assembly audit
(\texttt{AUDIT-\allowbreak{}F0-\allowbreak{}ASSEMBLY.\allowbreak{}md} \S{}1.1: the acronym UCP appeared in proposed
contracts with no canonical anchor in the definitions layer). Consensus
adoption of the universal textbook notion --- no project-specific content.
Consumers: the F0 lift rows (\texttt{lem-\allowbreak{}routef-\allowbreak{}f0-\allowbreak{}ucp-\allowbreak{}lift}), the K-ledger
factorization rows (UCP $\Delta,\Upsilon$), F2. Related:
\hyperref[def:stochastic]{\texttt{def-\allowbreak{}stochastic}} (row-stochastic = unital positive on $\ell_\infty^n$),
\hyperref[def:fd-cstar-diagonal]{\texttt{def-\allowbreak{}fd-\allowbreak{}cstar-\allowbreak{}diagonal}}.
\end{defnote}


